% Options for packages loaded elsewhere
\PassOptionsToPackage{unicode}{hyperref}
\PassOptionsToPackage{hyphens}{url}
\documentclass[
]{article}
\usepackage{xcolor}
\usepackage[margin=1in]{geometry}
\usepackage{amsmath,amssymb}
\setcounter{secnumdepth}{5}
\usepackage{iftex}
\ifPDFTeX
  \usepackage[T1]{fontenc}
  \usepackage[utf8]{inputenc}
  \usepackage{textcomp} % provide euro and other symbols
\else % if luatex or xetex
  \usepackage{unicode-math} % this also loads fontspec
  \defaultfontfeatures{Scale=MatchLowercase}
  \defaultfontfeatures[\rmfamily]{Ligatures=TeX,Scale=1}
\fi
\usepackage{lmodern}
\ifPDFTeX\else
  % xetex/luatex font selection
\fi
% Use upquote if available, for straight quotes in verbatim environments
\IfFileExists{upquote.sty}{\usepackage{upquote}}{}
\IfFileExists{microtype.sty}{% use microtype if available
  \usepackage[]{microtype}
  \UseMicrotypeSet[protrusion]{basicmath} % disable protrusion for tt fonts
}{}
\makeatletter
\@ifundefined{KOMAClassName}{% if non-KOMA class
  \IfFileExists{parskip.sty}{%
    \usepackage{parskip}
  }{% else
    \setlength{\parindent}{0pt}
    \setlength{\parskip}{6pt plus 2pt minus 1pt}}
}{% if KOMA class
  \KOMAoptions{parskip=half}}
\makeatother
\usepackage{color}
\usepackage{fancyvrb}
\newcommand{\VerbBar}{|}
\newcommand{\VERB}{\Verb[commandchars=\\\{\}]}
\DefineVerbatimEnvironment{Highlighting}{Verbatim}{commandchars=\\\{\}}
% Add ',fontsize=\small' for more characters per line
\newenvironment{Shaded}{}{}
\newcommand{\AlertTok}[1]{\textcolor[rgb]{1.00,0.00,0.00}{\textbf{#1}}}
\newcommand{\AnnotationTok}[1]{\textcolor[rgb]{0.38,0.63,0.69}{\textbf{\textit{#1}}}}
\newcommand{\AttributeTok}[1]{\textcolor[rgb]{0.49,0.56,0.16}{#1}}
\newcommand{\BaseNTok}[1]{\textcolor[rgb]{0.25,0.63,0.44}{#1}}
\newcommand{\BuiltInTok}[1]{\textcolor[rgb]{0.00,0.50,0.00}{#1}}
\newcommand{\CharTok}[1]{\textcolor[rgb]{0.25,0.44,0.63}{#1}}
\newcommand{\CommentTok}[1]{\textcolor[rgb]{0.38,0.63,0.69}{\textit{#1}}}
\newcommand{\CommentVarTok}[1]{\textcolor[rgb]{0.38,0.63,0.69}{\textbf{\textit{#1}}}}
\newcommand{\ConstantTok}[1]{\textcolor[rgb]{0.53,0.00,0.00}{#1}}
\newcommand{\ControlFlowTok}[1]{\textcolor[rgb]{0.00,0.44,0.13}{\textbf{#1}}}
\newcommand{\DataTypeTok}[1]{\textcolor[rgb]{0.56,0.13,0.00}{#1}}
\newcommand{\DecValTok}[1]{\textcolor[rgb]{0.25,0.63,0.44}{#1}}
\newcommand{\DocumentationTok}[1]{\textcolor[rgb]{0.73,0.13,0.13}{\textit{#1}}}
\newcommand{\ErrorTok}[1]{\textcolor[rgb]{1.00,0.00,0.00}{\textbf{#1}}}
\newcommand{\ExtensionTok}[1]{#1}
\newcommand{\FloatTok}[1]{\textcolor[rgb]{0.25,0.63,0.44}{#1}}
\newcommand{\FunctionTok}[1]{\textcolor[rgb]{0.02,0.16,0.49}{#1}}
\newcommand{\ImportTok}[1]{\textcolor[rgb]{0.00,0.50,0.00}{\textbf{#1}}}
\newcommand{\InformationTok}[1]{\textcolor[rgb]{0.38,0.63,0.69}{\textbf{\textit{#1}}}}
\newcommand{\KeywordTok}[1]{\textcolor[rgb]{0.00,0.44,0.13}{\textbf{#1}}}
\newcommand{\NormalTok}[1]{#1}
\newcommand{\OperatorTok}[1]{\textcolor[rgb]{0.40,0.40,0.40}{#1}}
\newcommand{\OtherTok}[1]{\textcolor[rgb]{0.00,0.44,0.13}{#1}}
\newcommand{\PreprocessorTok}[1]{\textcolor[rgb]{0.74,0.48,0.00}{#1}}
\newcommand{\RegionMarkerTok}[1]{#1}
\newcommand{\SpecialCharTok}[1]{\textcolor[rgb]{0.25,0.44,0.63}{#1}}
\newcommand{\SpecialStringTok}[1]{\textcolor[rgb]{0.73,0.40,0.53}{#1}}
\newcommand{\StringTok}[1]{\textcolor[rgb]{0.25,0.44,0.63}{#1}}
\newcommand{\VariableTok}[1]{\textcolor[rgb]{0.10,0.09,0.49}{#1}}
\newcommand{\VerbatimStringTok}[1]{\textcolor[rgb]{0.25,0.44,0.63}{#1}}
\newcommand{\WarningTok}[1]{\textcolor[rgb]{0.38,0.63,0.69}{\textbf{\textit{#1}}}}
\usepackage{longtable,booktabs,array}
\newcounter{none} % for unnumbered tables
\usepackage{calc} % for calculating minipage widths
% Correct order of tables after \paragraph or \subparagraph
\usepackage{etoolbox}
\makeatletter
\patchcmd\longtable{\par}{\if@noskipsec\mbox{}\fi\par}{}{}
\makeatother
% Allow footnotes in longtable head/foot
\IfFileExists{footnotehyper.sty}{\usepackage{footnotehyper}}{\usepackage{footnote}}
\makesavenoteenv{longtable}
\setlength{\emergencystretch}{3em} % prevent overfull lines
\providecommand{\tightlist}{%
  \setlength{\itemsep}{0pt}\setlength{\parskip}{0pt}}
\usepackage[]{natbib}
\bibliographystyle{plainnat}
% ============================================================================
% REQUIRED PACKAGES - Essential for document rendering
% ============================================================================

% Mathematical typesetting (required for equations and symbols)
\usepackage{amsmath,amssymb}          % Mathematical symbols and environments
\usepackage{amsfonts}                 % Additional math fonts
\usepackage{amsthm}                   % Theorem environments

% Graphics and page layout (required for figures and formatting)
\usepackage{graphicx}                 % Include graphics (REQUIRED for figures)
\usepackage[margin=1in]{geometry}     % Page margins
\usepackage{float}                    % Better float placement

% Tables (required for table formatting)
\usepackage{booktabs}                 % Professional tables
\usepackage{longtable}                % Long tables spanning pages
\usepackage{array}                    % Advanced table formatting

% PDF features (required for cross-references and metadata)
\usepackage{url}                      % URL formatting
\usepackage{hyperref}                 % Hyperlinks and cross-references
\usepackage{natbib}                   % Bibliography support (REQUIRED)

% ============================================================================
% PACKAGES - Improve formatting and functionality
% ============================================================================

% Table enhancements (optional but recommended)
\usepackage{multirow}                 % Multi-row table cells
\usepackage{caption}                  % caption formatting
\usepackage{subcaption}               % Sub-figures and sub-tables

% Math enhancements (optional but recommended)
\usepackage{bm}                       % Bold math symbols

% Reference enhancements (optional but recommended)
\usepackage{cleveref}                 % Intelligent cross-referencing
\usepackage{doi}                      % DOI links

% Configure figure numbering and captions
\renewcommand{\figurename}{Figure}
\captionsetup{
    justification=centering,
    font=small,
    labelfont=bf,
    labelsep=period
}

% Configure table numbering and captions
\renewcommand{\tablename}{Table}
\captionsetup[table]{
    justification=centering,
    font=small,
    labelfont=bf,
    labelsep=period
}

% Configure section numbering
\setcounter{secnumdepth}{3}
\renewcommand{\thesection}{\arabic{section}}
\renewcommand{\thesubsection}{\arabic{section}.\arabic{subsection}}
\renewcommand{\thesubsubsection}{\arabic{section}.\arabic{subsection}.\arabic{subsubsection}}

% Configure equation numbering
\numberwithin{equation}{section}

% Configure hyperref for proper linking
\hypersetup{
    colorlinks=true,
    linkcolor=red,
    citecolor=red,
    urlcolor=red,
    filecolor=red,
    pdfborder={0 0 0},
    bookmarks=true,
    bookmarksnumbered=true,
    bookmarkstype=toc,
    pdftitle={Ento-Linguistic Domains: Language, Ambiguity, and Scientific Communication in Entomology},
    pdfauthor={Daniel Ari Friedman},
    pdfsubject={Scientific Terminology Analysis in Entomology},
    pdfkeywords={entomology, scientific terminology, discourse analysis, anthropomorphism, language, ant biology},
    pdfcreator={render_pdf.sh},
    pdfproducer={XeLaTeX}
}

% ============================================================================
% PACKAGE CONFIGURATION
% ============================================================================

% Configure cleveref for intelligent cross-references
\crefname{section}{Section}{Sections}
\crefname{subsection}{Subsection}{Subsections}
\crefname{subsubsection}{Subsubsection}{Subsubsections}
\crefname{equation}{Equation}{Equations}
\crefname{figure}{Figure}{Figures}
\crefname{table}{Table}{Tables}
\crefname{appendix}{Appendix}{Appendices}

% Configure fonts for Unicode support with fallbacks
\usepackage{newunicodechar}
%\newunicodechar definitions for Unicode characters used in the manuscript
\newunicodechar{⁴}{\textsuperscript{4}}
\newunicodechar{₄}{\textsubscript{4}}
\newunicodechar{²}{\textsuperscript{2}}
\newunicodechar{₀}{\textsubscript{0}}
\newunicodechar{₁}{\textsubscript{1}}
\newunicodechar{₂}{\textsubscript{2}}
\newunicodechar{₃}{\textsubscript{3}}
% Arrow characters used in abstract, conclusion, and supplemental methods
\newunicodechar{→}{\ensuremath{\rightarrow}}
\newunicodechar{←}{\ensuremath{\leftarrow}}
\newunicodechar{↔}{\ensuremath{\leftrightarrow}}
% Additional symbols
\newunicodechar{≥}{\ensuremath{\geq}}
\newunicodechar{≤}{\ensuremath{\leq}}
\newunicodechar{≠}{\ensuremath{\neq}}
\newunicodechar{…}{\ldots}

% ============================================================================
% FONTS AND TYPOGRAPHY
% ============================================================================

% Use standard fonts for better compatibility
\usepackage{lmodern}
\usepackage[T1]{fontenc}

% ============================================================================
% CODE BLOCK STYLING
% ============================================================================

% code block styling for better contrast and readability
\usepackage{fancyvrb}
\usepackage{fvextra}   % Extends fancyvrb: adds bgcolor, breaklines, etc.
\usepackage{xcolor}
\usepackage{listings}

% Define custom colors for code blocks
\definecolor{codebg}{RGB}{248, 248, 248}      % Very light gray background
\definecolor{codeborder}{RGB}{200, 200, 200}  % Medium gray border
\definecolor{codefg}{RGB}{34, 34, 34}         % Dark gray text
\definecolor{commentcolor}{RGB}{102, 102, 102} % Comment color
\definecolor{keywordcolor}{RGB}{0, 0, 0}       % Keyword color
\definecolor{stringcolor}{RGB}{0, 102, 0}      % String color

% Configure Verbatim environment for inline code
\DefineVerbatimEnvironment{Verbatim}{Verbatim}{%
    fontsize=\small,
    frame=single,
    framerule=0.5pt,
    framesep=3pt,
    rulecolor=\color{codeborder},
    bgcolor=codebg
}

% Configure code block styling
\DefineVerbatimEnvironment{Highlighting}{Verbatim}{%
    fontsize=\footnotesize,
    frame=single,
    framerule=0.5pt,
    framesep=5pt,
    rulecolor=\color{codeborder},
    bgcolor=codebg
}

% Style inline code with \texttt
\renewcommand{\texttt}[1]{%
    \colorbox{codebg}{\color{codefg}\ttfamily #1}%
}

% Configure listings package for code blocks
\lstset{
    backgroundcolor=\color{codebg},
    basicstyle=\footnotesize\ttfamily\color{codefg},
    breakatwhitespace=false,
    breaklines=true,
    captionpos=b,
    commentstyle=\color{commentcolor},
    deletekeywords={...},
    escapeinside={\%*}{*)},
    extendedchars=true,
    frame=single,
    framerule=0.5pt,
    framesep=5pt,
    keepspaces=true,
    keywordstyle=\color{keywordcolor}\bfseries,
    language=Python,
    morekeywords={*,...},
    numbers=left,
    numbersep=5pt,
    numberstyle=\tiny\color{codefg},
    rulecolor=\color{codeborder},
    showspaces=false,
    showstringspaces=false,
    showtabs=false,
    stepnumber=1,
    stringstyle=\color{stringcolor},
    tabsize=4,
    title=\lstname
}

% Override any Pandoc default lstset configurations
\AtBeginDocument{
    \lstset{
        backgroundcolor=\color{codebg},
        basicstyle=\footnotesize\ttfamily\color{codefg},
        frame=single,
        framerule=0.5pt,
        framesep=5pt,
        rulecolor=\color{codeborder},
        numbers=left,
        numbersep=5pt,
        numberstyle=\tiny\color{codefg}
    }
}

% Configure bibliography
% Note: Using plainnat with natbib package for proper citation processing
% The bibliography style and commands (\bibliographystyle and \bibliography) are in 99_references.md

% Simple page break support for document structure
% Note: Page breaks are handled in the markdown generation, not here

% ============================================================================
% DOCUMENT FORMATTING
% ============================================================================

% Ensure proper spacing and formatting
\frenchspacing  % Single space after periods
\linespread{1.2}  % Slightly increased line spacing for readability
\usepackage{bookmark}
\IfFileExists{xurl.sty}{\usepackage{xurl}}{} % add URL line breaks if available
\urlstyle{same}
% Make links footnotes instead of hotlinks:
\DeclareRobustCommand{\href}[2]{#2\footnote{\url{#1}}}
\hypersetup{
  pdftitle={Ento-Linguistic Domains: Language{,} Ambiguity{,} and Scientific Communication in Entomology},
  pdfauthor={Daniel Ari Friedman},
  hidelinks,
  pdfcreator={LaTeX via pandoc}}

\title{Ento-Linguistic Domains: Language, Ambiguity, and Scientific
Communication in Entomology}
\usepackage{etoolbox}
\makeatletter
\providecommand{\subtitle}[1]{% add subtitle to \maketitle
  \apptocmd{\@title}{\par {\large #1 \par}}{}{}
}
\makeatother
\subtitle{How Terminology Networks Shape Understanding of Insect Biology
(And Vice-Versa)}
\author{Daniel Ari Friedman}
\date{February 24, 2026}

\begin{document}
\maketitle

{
\setcounter{tocdepth}{2}
\tableofcontents
}
\section{Abstract}\label{sec:abstract}

Scientific language does not merely describe biological phenomena; it
constitutes the \textbf{generative model} through which researchers
parse complex systems. This paper makes three contributions to
understanding---and correcting---the epistemic consequences of that
constitutive role. First, we introduce a \textbf{six-domain
Ento-Linguistic framework} that decomposes the terminological landscape
of ant research into analytically tractable categories where
anthropomorphic language most distorts causal modeling. Second, we
develop an \textbf{open-source computational pipeline} integrating
automated term extraction, co-occurrence network construction, and
information-theoretic ambiguity scoring with \textbf{Active Inference}
and \textbf{Complex Systems Theory}. Third, we propose and validate four
evidence-based meta-standards---\textbf{Clarity},
\textbf{Appropriateness}, \textbf{Consistency}, and
\textbf{Evolvability} (CACE)---as a protocol for \textbf{lexical
engineering}.

Analysis of a corpus of \textbf{362 entomological publications} (48,062
tokens; 6,971 unique token types) extracts \textbf{871 candidate terms}
(223 assigned to specific domains) across six Ento-Linguistic domains,
revealing terminology networks with strong modularity and systematic
cross-domain bridging. The Power \& Labor domain consistently yields the
highest ambiguity scores (0.820): hub terms such as ``caste'' and
``queen'' impose hierarchical control layers absent from the stigmergic
biology they purport to describe. Each domain's top terms emerge
directly from the literature: corpus-level frequency leaders include
\emph{ant} (1,032), \emph{colony} (838), \emph{worker} (813),
\emph{queen} (592), \emph{social} (586), and \emph{insect} (297). Across
all 223 domain-assigned terms, 73.4\% exhibit context-dependent
meanings, and conceptual networks reveal that ``individuality'' spans
multiple biological scales, blurring the formal boundaries
(\textbf{Markov Blankets}) required for rigorous systems modeling.

CACE validation on the historical ``slave'' → ``host worker'' reform
confirms the framework's diagnostic and prescriptive power: aggregate
scores rise from 0.38 to 0.81, with the largest gain in Appropriateness
(+0.60). These findings extend beyond entomology to any field where
human social concepts are projected onto non-human systems. The
accompanying reproducible pipeline provides the analytical tools for a
more self-aware and epistemically rigorous scientific practice.

\newpage

\section{Introduction}\label{sec:introduction}

\subsection{Linguistic Priors and Generative
Models}\label{linguistic-priors-and-generative-models}

Scientific inquiry is a process of \textbf{active inference}, where
researchers refine generative models to minimize surprise about
biological observations \cite{friston2010free}. Language acts as the
\textbf{hyper-prior} for these models: it constrains the hypothesis
space before data collection even begins. When entomologists employ
terms like ``queen'' or ``caste,'' they are not merely labeling
phenomena; they are importing a high-precision prior from human social
systems into their model of insect biology. If this prior is
structurally misaligned with the target system---for instance, assuming
top-down control in a stigmergic network---the resulting model will
necessarily suffer from high variational free energy, manifesting as
persistent anomalies and ``epicycles'' in theoretical explanations
\cite{kuhn1996, clark2013whatever}.

We can formally model the \textbf{scientific community itself as a
multi-scale Active Inference agent}. Its collective task is to minimize
long-term surprise (variational free energy) about the entomological
world it observes. Its ``Generative Model'' is the shared ontology of
the field---the lexicon and conceptual structures encoded in the
literature. When this ontology is precise and plastic, the community
efficiently updates its priors in response to new sensory data (e.g.,
genomic evidence). However, when the ontology is rigid or laden with
hidden anthropomorphic priors, the agent suffers from \textbf{Prior
Dogmatism}: a failure of belief updating where high-precision, fixed
priors overwhelm contradictory sensory evidence. In this state,
anomalies are explained away rather than used to update the model.
Terminology reform is thus not a political act but a \textbf{model
selection} process: optimizing the community's generative model to
restore its ability to minimize free energy.

This optimization requires specific criteria. We propose
\textbf{Evolvability}---defined here as \textbf{scale-invariance}---as a
critical metric for scientific terms. An evolvable term maintains its
validity across biological scales (gene, organism, superorganism)
without fracturing. A term like ``queen,'' by contrast, is
scale-brittle: it functions as a metaphor at the colony level but
dissolves into incoherence when applied to the underlying genetic or
molecular mechanisms of caste determination.

Recent formal work by \citet{friedman2021active} demonstrates that ant
colonies can be modeled as ensembles of \emph{active
inferants}---individual agents performing Bayesian inference over local
states via chemical stigmergy---without any centralized controller; yet
the dominant vocabulary of the field continues to presuppose one.

Our work examines this epistemic risk through systematic analysis of
\emph{Ento-Linguistic domains}: specific areas where linguistic priors
obscure the causal architecture of ant systems.

The consequences of this misalignment are not merely philosophical. They
propagate through every stage of the research cycle---from the
hypotheses a researcher formulates, through the variables selected for
measurement, to the causal explanations offered for observed phenomena.
The following section formalizes this propagation as a problem of model
integrity.

\subsection{Motivation: Minimizing Model
Misspecification}\label{motivation-minimizing-model-misspecification}

The drive for clarity is not merely a stylistic preference but a
requirement for model integrity. As \citet{keller1995language} argued,
the language of science constitutes the cognitive scaffolding of
research. In the framework of Active Inference, an undefined or
metaphor-laden term introduces \textbf{irreducible uncertainty}
(entropy) into the scientific communication channel.

The present moment demands this formalization. Recent cognitive science
emphasizes that metaphor is a mechanism of predictive processing
\cite{steen2017deliberate}. Rather than perpetuating ``legacy code'' in
our linguistic ontology, researchers must critically assess whether
their terminological priors minimize or maximize the complexity of their
biological models.

A paradigmatic example is the ``slave-making'' debate.
\citet{herbers2006} showed that the term ``slave'' naturalizes a human
institution while obscuring the biological mechanism of \textbf{social
parasitism}. In formal terms, the ``slave'' metaphor implies a conscious
coercion policy, whereas the replacement term ``dulosis'' correctly
identifies the phenomenon as a breakdown in nestmate recognition signals
(a failure of the Markov Blanket's security filter). Reform, therefore,
is not just about ethics; it is about restoring the causal fidelity of
the scientific model.

\subsection{The Challenge of Terminological
Reform}\label{the-challenge-of-terminological-reform}

A common objection to improving scientific language is that changing
terminology creates disconnection from existing literature. If
entomologists abandon terms like ``caste'' or ``slave,'' how would
researchers locate papers on task performance or social parasitism?

This objection, however, inadvertently strengthens the case for reform.
Retaining problematic terminology for convenience perpetuates and
compounds conceptual distortions rather than addressing them
\cite{herbers2006}. The appropriate response is to work systematically
toward clearer communication while developing the necessary
infrastructure for literature synthesis---restructuring information from
existing sources and establishing new meta-standards for scientific
discourse. Recent community-level momentum confirms this trajectory:
discussions at the MirMeco 2023 International Ant Meeting
\cite{laciny2024terminology} and the Entomological Society of America's
Better Common Names Project \cite{betternamesproject2024} demonstrate
that the professional community increasingly shares these concerns.

\subsection{Ento-Linguistic Domains: A Framework for
Analysis}\label{ento-linguistic-domains-a-framework-for-analysis}

We organize our analysis around six domains where entomological language
creates ambiguity or imports unjustified assumptions. Each domain
isolates a distinct category of terminological friction between human
conceptual frameworks and ant biology.

\textbf{Unit of Individuality.} The definition of a biological
individual is formally equivalent to the specification of a
\textbf{Markov Blanket}---the statistical boundary separating internal
states from external states \cite{friston2013life}. Terms like
``colony,'' ``superorganism,'' and ``individual'' confuse these
boundaries, creating models where the relevant unit of agency is
undefined.

\textbf{Behavior and Identity.} Task performance in ants is a fluid
process of \textbf{policy selection} based on local cues
\cite{gordon2010}. However, terminology transforms these transient
policies into categorical identities (``forager,'' ``nurse''). This
effectively hard-codes a fixed-role prior into the model, obscuring the
plasticity and Bayesian updating that actually drives task allocation.

\textbf{Power \& Labor.} Terms like ``queen,'' ``worker,'' and ``caste''
impose a hierarchical control architecture on a system that is
fundamentally \textbf{stigmergic}. This introduces a causal error: it
attributes colony-level regulation to centralized agency (the queen)
rather than distributed feedback loops, fundamentally misrepresenting
the system's control theory.

\textbf{Sex \& Reproduction.} Terms like ``sex determination'' and ``sex
differentiation'' carry implicit assumptions about binary systems that
may not map onto ant reproductive biology, where haplodiploidy creates
fundamentally different patterns \cite{chandra2021epigenetics}.

\textbf{Kin \& Relatedness.} Human kinship terminology, grounded in
bilateral relatedness, creates systematic friction when applied to ant
societies structured by haplodiploidy. In haplodiploid species, full
sisters share an average relatedness coefficient of
\(r = 0.75\)---higher than the mother--daughter coefficient of
\(r = 0.5\)---a fundamental asymmetry absent from human kinship models.
Terms such as ``sister,'' ``mother,'' and ``family'' obscure this
asymmetry and its profound consequences for kin selection theory
\cite{chandra2021epigenetics}.

\textbf{Economics.} Economic metaphors---markets, trade, investment,
cost-benefit---shape analysis of ant foraging, resource distribution,
and colony-level resource management. This domain investigates how
transactional frameworks constrain biological interpretation by
conflating proximate energetic expenditure with ultimate fitness costs,
importing assumptions of rational optimisation from microeconomics into
systems that operate through evolved heuristics rather than deliberative
calculation. In Active Inference terms, economic metaphors impose a
\textbf{utility-maximising} generative model on systems that instead
minimise variational free energy through local policy selection---a
distinction with profound consequences for how foraging efficiency,
brood investment, and inter-colony resource flows are modelled and
interpreted.

\subsection{Research Approach}\label{research-approach}

This work employs a mixed-methodology framework combining computational
text analysis with theoretical discourse examination. The computational
component processes a \textbf{corpus of 362 entomological publications}
(48,062 tokens; 6,971 unique token types; 871 extracted candidate terms,
223 domain-assigned) using automated term extraction, co-occurrence
network construction, and ambiguity scoring to identify statistical
patterns in language use. The theoretical component, informed by
\citeauthor{foucault1972archaeology}'s archaeological method
\citeyearpar{foucault1972archaeology}, conceptual metaphor theory
\cite{lakoff1980metaphors}, and \citeauthor{gordon2023ecology}'s
\citeyearpar{gordon2023ecology} ecological framework for collective
behaviour, examines how these patterns reflect deeper conceptual
structures. Longitudinal case studies of the ``caste'' and
``superorganism'' vocabularies (Section \ref{sec:experimental_results})
track how terminology has evolved alongside empirical discoveries over
five decades, providing diachronic evidence for the framework's claims.
All data and analysis code are reproducible and available for validation
and extension.

\subsection{Terminology Network
Visualization}\label{terminology-network-visualization}

To illustrate the framework's output, Figure \ref{fig:concept_map} shows
how terms cluster around the six Ento-Linguistic domains and form
cross-domain networks of meaning; detailed quantitative analysis follows
in Section \ref{sec:experimental_results}.

\begin{figure}[h]
\centering
\includegraphics[width=0.9\textwidth]{/Users/4d/Documents/GitHub/template/projects/ento_linguistics/output/figures/concept_map.png}
\caption{Conceptual map of Ento-Linguistic domains showing relationships between terminology networks. Each node represents an extracted concept; node size is proportional to term frequency in the corpus and node colour encodes the primary domain assignment. Edges connect co-occurring concepts, with thickness reflecting co-occurrence strength. The six domains form interconnected clusters; central hub terms such as ``colony,''``caste,'' and ``individual'' bridge multiple domains, demonstrating how specific terminological choices propagate across the scientific discourse of entomology.}
\label{fig:concept_map}
\end{figure}

\newpage

\section{Methodology}\label{sec:methodology}

Our research employs a mixed-methodology framework combining two
sequential phases: first, systematic corpus construction through
multi-source literature mining; second, a multi-layer computational and
discourse-analytic pipeline applied to the assembled corpus. This
section describes each phase at the level of detail required for
replication; full implementation specifications are provided in
Supplemental Methods \ref{sec:supplemental_methods}, and extended
theoretical and analytical derivations in Supplemental Analysis
\ref{sec:supplemental_analysis}.

\begin{center}\rule{0.5\linewidth}{0.5pt}\end{center}

\subsection{3.1 Data Acquisition}\label{sec:data_acquisition}

\subsubsection{3.1.1 Search Strategy and Source
Selection}\label{search-strategy-and-source-selection}

The primary corpus was assembled by querying two open-access databases
--- \textbf{PubMed} (NCBI) and \textbf{arXiv} --- using the
\texttt{PubMedMiner} and \texttt{ArXivMiner} classes implemented in
\texttt{src/data/literature\_mining.py}. Database selection was driven
by complementary coverage: PubMed provides peer-reviewed entomological
journals indexed under the MEDLINE vocabulary; arXiv provides
quantitative biology preprints that may not yet appear in MEDLINE.

\textbf{PubMed query} (\texttt{create\_entomology\_query()}):

\begin{Shaded}
\begin{Highlighting}[]
\NormalTok{"ants" OR "Formicidae" OR "Hymenoptera" OR "eusocial" OR "eusociality"}
\NormalTok{OR "social insects" OR "colony" OR "nest" OR "foraging"}
\NormalTok{OR "division of labor" AND (English[Language])}
\end{Highlighting}
\end{Shaded}

The query was submitted to the NCBI Entrez Utilities API
(\texttt{esearch.fcgi}, \texttt{retmode=json}). Results were retrieved
in batches of 20 PMIDs via \texttt{eSummary} (metadata: title, authors,
journal, year, DOI) and \texttt{eFetch} (\texttt{rettype=abstract},
\texttt{retmode=xml}) for abstract text, with a 500 ms inter-batch delay
to comply with NCBI rate limits. \texttt{PubMedMiner} caches search
results (\texttt{enable\_cache=True}) to prevent redundant API calls
during pipeline re-runs.

\textbf{arXiv query} (\texttt{ArXivMiner.search()}):

\begin{Shaded}
\begin{Highlighting}[]
\NormalTok{cat:q{-}bio.PE OR cat:q{-}bio.QM}
\end{Highlighting}
\end{Shaded}

Results retrieved via the \href{http://export.arxiv.org/api/query}{arXiv
API} (Atom XML, \texttt{sortBy=submittedDate},
\texttt{sortOrder=descending}). Records were post-filtered by keyword
overlap against the entomology vocabulary set \{\emph{ant, ants,
formicidae, eusocial, colony, social insect}\}; only records whose
combined title+abstract text contained at least one keyword were
retained.

\subsubsection{3.1.2 Corpus Composition and
Cleaning}\label{corpus-composition-and-cleaning}

The assembled \texttt{LiteratureCorpus} stores \texttt{Publication}
dataclass objects with the following fields: \texttt{title},
\texttt{authors}, \texttt{abstract}, \texttt{doi}, \texttt{pmid},
\texttt{year}, \texttt{journal}, \texttt{keywords}, \texttt{full\_text}.
After deduplication (by PMID) and quality filtering (records with
neither abstract nor full text were excluded), the final corpus
comprises:

{\def\LTcaptype{none} % do not increment counter
\begin{longtable}[]{@{}ll@{}}
\toprule\noalign{}
Metric & Value \\
\midrule\noalign{}
\endhead
\bottomrule\noalign{}
\endlastfoot
Documents & \textbf{369} \\
Total processed tokens & \textbf{48,952} \\
Unique token types & \textbf{7,067} \\
Type--token ratio & \textbf{0.144} \\
Total characters & \textbf{539,570} \\
Average token length & \textbf{7.32 characters} \\
\end{longtable}
}

These statistics are computed at runtime by
\texttt{TextProcessor.get\_vocabulary\_stats()} and serialized to
\texttt{output/data/corpus\_statistics.json}; the values reported here
are read directly from that file and are therefore always current with
the last pipeline run.

Full text preprocessing --- tokenization (\texttt{sent\_tokenize},
\texttt{word\_tokenize}), scientific-term merge, stop-word removal (NLTK
English + \texttt{SCIENTIFIC\_STOP\_WORDS}), and lemmatization
(\texttt{WordNetLemmatizer}) --- is documented in Section S2 of
Supplemental Methods \ref{sec:supplemental_methods}.

\subsubsection{3.1.3 Domain Coverage
Verification}\label{domain-coverage-verification}

To verify that the search strategy captured all six target
Ento-Linguistic domains, term counts were computed across domains
immediately after corpus construction. The six domains and their seed
vocabularies are:

{\def\LTcaptype{none} % do not increment counter
\begin{longtable}[]{@{}
  >{\raggedright\arraybackslash}p{(\linewidth - 2\tabcolsep) * \real{0.2857}}
  >{\raggedright\arraybackslash}p{(\linewidth - 2\tabcolsep) * \real{0.7143}}@{}}
\toprule\noalign{}
\begin{minipage}[b]{\linewidth}\raggedright
Domain
\end{minipage} & \begin{minipage}[b]{\linewidth}\raggedright
Example Seed Terms
\end{minipage} \\
\midrule\noalign{}
\endhead
\bottomrule\noalign{}
\endlastfoot
Power \& Labor & queen, worker, dominance, hierarchy, division of
labor \\
Unit of Individuality & colony, superorganism, eusocial, individual,
organism \\
Sex \& Reproduction & mating, haplodiploidy, parthenogenesis, queen,
egg \\
Behavior \& Identity & caste, forager, nurse, task, polyethism \\
Kin \& Relatedness & kin selection, inclusive fitness, relatedness,
altruism \\
Economics & foraging, cost, benefit, resource allocation, trade-off \\
\end{longtable}
}

The current pipeline run extracted \textbf{871 terms distributed across
all six domains}, of which \textbf{223 receive specific domain
assignments}, sourced from \texttt{output/data/domain\_statistics.json}.
Domain-specific acquisition details, bridging term frequencies, and
per-domain confidence statistics are reported in Supplemental Results
\ref{sec:supplemental_results}.

\begin{center}\rule{0.5\linewidth}{0.5pt}\end{center}

\subsection{3.2 Statistical Analysis}\label{sec:statistical_analysis}

\subsubsection{3.2.1 Analytical Framework
Overview}\label{analytical-framework-overview}

The statistical pipeline comprises six interdependent analytical layers
applied sequentially to the assembled corpus: (1) term extraction and
classification, (2) semantic entropy estimation, (3) domain-level
statistical testing, (4) conceptual network construction and centrality
analysis, (5) rhetorical and discourse pattern scoring, and (6) CACE
meta-standard evaluation. All analyses are implemented in
\texttt{src/analysis/} and are fully deterministic
(\texttt{random\_state=42} throughout). Extended statistical derivations
and robustness tests are presented in Supplemental Analysis
\ref{sec:supplemental_analysis}.

\subsubsection{3.2.2 Term Extraction and
Classification}\label{term-extraction-and-classification}

\texttt{TerminologyExtractor}
(\texttt{src/analysis/term\_extraction.py}) assigns each extracted term
to one or more domains via seed-expansion: tokens are first matched
against a domain seed lexicon, then extended to co-occurring tokens
within a 5-token sliding window. Each \texttt{Term} dataclass carries
\texttt{text}, \texttt{lemma}, \texttt{domains}, \texttt{frequency},
\texttt{contexts} (deduplicated sentences), \texttt{pos\_tags},
\texttt{confidence}, and \texttt{semantic\_entropy}. N-gram extraction
(\texttt{TextProcessor.extract\_ngrams}) captures compound terms (e.g.,
\emph{division of labor}, \emph{kin selection}) that single-token
analysis would fragment. Full API documentation is in Section S3 of
Supplemental Methods \ref{sec:supplemental_methods}.

\subsubsection{3.2.3 Semantic Entropy}\label{semantic-entropy}

To quantify terminological ambiguity, we compute \textbf{Semantic
Entropy} \(H(t)\) for each term \(t\) with sufficient attestation
(\(\geq 5\) valid contexts):

\begin{equation}\label{eq:semantic_entropy}
H(t) = -\sum_{i=1}^{k} p_i \log_2 p_i \qquad \text{(bits)}
\end{equation}

where \(p_i\) is the empirical proportion of usage contexts assigned to
semantic cluster \(i\) by \(k\)-means (\(k = \min(5, |C_t|)\),
scikit-learn, \texttt{random\_state=42}) over TF-IDF context vectors.
Terms with \(H(t) > H^* = 2.0\) bits --- corresponding to four or more
equiprobable semantic senses --- are flagged \texttt{is\_high\_entropy}.
The threshold was calibrated against terms of known polysemy
(\emph{colony}, \emph{queen}) and specificity (\emph{haplodiploidy},
\emph{trophallaxis}). Implementation:
\texttt{src/analysis/semantic\_entropy.py::calculate\_semantic\_entropy};
corpus-level results:
\texttt{src/analysis/semantic\_entropy.py::calculate\_corpus\_entropy}.

\subsubsection{3.2.4 Domain-Level Statistical
Tests}\label{domain-level-statistical-tests}

Cross-domain entropy comparisons use the following battery, all
implemented from scratch in \texttt{src/analysis/statistics.py}:

{\def\LTcaptype{none} % do not increment counter
\begin{longtable}[]{@{}
  >{\raggedright\arraybackslash}p{(\linewidth - 4\tabcolsep) * \real{0.2069}}
  >{\raggedright\arraybackslash}p{(\linewidth - 4\tabcolsep) * \real{0.3448}}
  >{\raggedright\arraybackslash}p{(\linewidth - 4\tabcolsep) * \real{0.4483}}@{}}
\toprule\noalign{}
\begin{minipage}[b]{\linewidth}\raggedright
Test
\end{minipage} & \begin{minipage}[b]{\linewidth}\raggedright
Function
\end{minipage} & \begin{minipage}[b]{\linewidth}\raggedright
Application
\end{minipage} \\
\midrule\noalign{}
\endhead
\bottomrule\noalign{}
\endlastfoot
Two-sample Welch's \(t\)-test & \texttt{t\_test} & Pairwise entropy
comparison between domains \\
One-way ANOVA & \texttt{anova\_test} & Simultaneous entropy comparison
across all 6 domains \\
95\% confidence intervals & \texttt{calculate\_confidence\_interval} &
Mean entropy uncertainty per domain \\
Pearson / Spearman correlation & \texttt{calculate\_correlation} &
Entropy--frequency relationship \\
Normal / Exponential / Uniform fit & \texttt{fit\_distribution} &
Entropy distribution characterization \\
\end{longtable}
}

The Welch--Satterthwaite degrees-of-freedom approximation is applied in
all two-sample \(t\)-tests; \(p\)-values are computed via
\texttt{scipy.stats.t.sf} and \texttt{scipy.stats.f.sf}. For the
\(\binom{6}{2} = 15\) pairwise domain comparisons, \(p\)-values are
corrected using the \textbf{Benjamini--Hochberg} false discovery rate
procedure at \(q = 0.05\). Effect sizes are reported as Cohen's \(d\)
(small: 0.2, medium: 0.5, large: 0.8 \cite{cohen1988statistical}).

A three-level \textbf{multi-scale ambiguity classification} is applied
to high-entropy terms: (1) \emph{Lexical Ambiguity} --- multiple
dictionary meanings; (2) \emph{Scale Ambiguity} --- meaning variation
across biological scales
(\(\text{gene} \to \text{organism} \to \text{colony}\)); (3)
\emph{Temporal Ambiguity} --- historical meaning evolution traceable
across publication years. The biological-scale dimension is further
formalized through the \textbf{Markov Blanket} formalism
\cite{friston2013life}.

\subsubsection{3.2.5 Conceptual Network
Analysis}\label{conceptual-network-analysis}

\texttt{ConceptualMapper} (\texttt{src/analysis/conceptual\_mapping.py})
constructs a \texttt{ConceptMap} of \textbf{6 concepts}
(biological\_individuality, social\_organization, reproductive\_biology,
kinship\_systems, resource\_economics, behavioral\_ecology) linked by
\textbf{8 weighted edges}. Edge weights are Jaccard overlap ratios:

\begin{equation}\label{eq:jaccard_overlap}
w_{AB} = \frac{|A \cap B|}{\min(|A|, |B|)}
\end{equation}

Composite relationship strength decomposes as:
\(\text{strength} = 0.4\,w_\text{base} + 0.3\,r_\text{term} + 0.2\,r_\text{domain} + 0.1\,\mathbb{1}_\text{hierarchical}\).

Centrality analysis uses NetworkX: degree centrality, betweenness
centrality, closeness centrality, and eigenvector centrality
(\texttt{max\_iter=1000}; \texttt{PowerIterationFailedConvergence}
fallback \(\to\) 0). Concept-level results are serialized to
\texttt{output/data/concept\_map\_summary.json}. Cross-domain bridging
terms --- appearing in \(\geq 2\) domains --- are identified with
\texttt{identify\_cross\_domain\_bridges}; the current run yields 43
bridging terms in Power \& Labor and 26 in Sex \& Reproduction.

\subsubsection{3.2.6 Rhetorical and Discourse
Analysis}\label{rhetorical-and-discourse-analysis}

\texttt{analyze\_rhetorical\_strategies}
(\texttt{src/analysis/rhetorical\_analysis.py}) detects four strategy
types per abstract via regex:

{\def\LTcaptype{none} % do not increment counter
\begin{longtable}[]{@{}
  >{\raggedright\arraybackslash}p{(\linewidth - 2\tabcolsep) * \real{0.4762}}
  >{\raggedright\arraybackslash}p{(\linewidth - 2\tabcolsep) * \real{0.5238}}@{}}
\toprule\noalign{}
\begin{minipage}[b]{\linewidth}\raggedright
Strategy
\end{minipage} & \begin{minipage}[b]{\linewidth}\raggedright
Detection
\end{minipage} \\
\midrule\noalign{}
\endhead
\bottomrule\noalign{}
\endlastfoot
Authority &
\texttt{\textbackslash{}(.*?20\textbackslash{}d\{2\}.*?\textbackslash{})}
--- citation parentheticals \\
Analogy &
\texttt{\textbackslash{}blike\textbackslash{}s+.*?\textbackslash{}bant}
--- ant-comparison expressions \\
Generalization &
\texttt{\textbackslash{}b(all\textbackslash{}\textbar{}every\textbackslash{}\textbar{}always\textbackslash{}\textbar{}never)\textbackslash{}s+.*?\textbackslash{}bant}
--- absolutist quantifiers \\
Anecdotal &
\texttt{\textbackslash{}b(for\ example\textbackslash{}\textbar{}such\ as\textbackslash{}\textbar{}consider\textbackslash{}\textbar{}imagine)\textbackslash{}b}
--- evidential markers \\
\end{longtable}
}

\texttt{identify\_narrative\_frameworks} classifies each abstract into
one or more of four narrative types (progress, conflict, discovery,
complexity) by keyword presence.
\texttt{score\_argumentative\_structures} decomposes argumentative
strength into claim strength, evidence quality, and reasoning coherence,
averaged to an overall score. \texttt{LinguisticFeatureExtractor}
computes anthropomorphic (\(4\) patterns), hierarchical (\(4\)
patterns), and economic (\(4\) patterns) framing densities per document.
Anthropomorphic framing indicators include five conceptual categories
--- agency, communication, social contract, cognition, and hierarchy ---
as specified in
\texttt{ConceptualMapper.detect\_anthropomorphic\_concepts()}.

\subsubsection{3.2.7 CACE Evaluation}\label{cace-evaluation}

Each term is scored on four bounded \([0,1]\) dimensions constituting
the \textbf{CACE} framework (\texttt{src/analysis/cace\_scoring.py}):

\begin{equation}\label{eq:cace_clarity}
\text{Clarity}(t) = \max\!\left(0,\, 1 - \frac{H(t)}{\log_2 5}\right)
\end{equation}

\begin{equation}\label{eq:cace_appropriateness}
\text{Appropriateness}(t) = 0 \text{ if } t \in \mathcal{A},\; \text{else } \propto \text{context-match ratio}
\end{equation}

\begin{equation}\label{eq:cace_consistency}
\text{Consistency}(t) = 1 - \operatorname{Var}\!\left(\mathbf{X}_t\right),\; \mathbf{X}_t = \text{TF-IDF context matrix}
\end{equation}

\begin{equation}\label{eq:cace_evolvability}
\text{Evolvability}(t) = \frac{|\text{SCALE\_LEVELS represented in contexts}|}{3}
\end{equation}

where \(\mathcal{A}\) is the \texttt{ANTHROPOMORPHIC\_TERMS} set
(\(\ni\) queen, king, slave, worker, soldier, nurse, \ldots). The
aggregate CACE score is the arithmetic mean of the four dimensions.
\texttt{compare\_terms\_cace} returns a ranked list for all terms.
Inter-rater reliability for qualitative CACE audits is assessed via
Cohen's \(\kappa\). Full formulas and the \texttt{CACEScore} dataclass
specification are in \texttt{src/analysis/cace\_scoring.py}.

\subsubsection{3.2.8 Validation and
Reproducibility}\label{validation-and-reproducibility}

Results are validated through three mechanisms: (1) \textbf{internal
consistency} --- term frequency distributions against semantic entropy
estimates; (2) \textbf{cross-method agreement} --- rhetorical pattern
frequencies against domain framing scores; (3) \textbf{external
triangulation} --- comparison against existing critical discourse
analyses of entomological literature \cite{latour1987, longino1990}.
Robustness testing (subsampling stability, parameter sensitivity,
annotation consistency) is implemented throughout
\texttt{src/analysis/}.

The pipeline is fully deterministic and clean-slate:
\texttt{output/figures/} and \texttt{output/data/} are wiped and
recreated on every run
(\texttt{scripts/02\_generate\_figures.py::\_setup\_directories}),
ensuring no stale artefacts persist across runs. All corpus statistics
cited in §3.1 are read at runtime from
\texttt{output/data/corpus\_statistics.json} and are never hardcoded.

\newpage

\section{Experimental Results}\label{sec:experimental_results}

\subsection{Terminology Extraction Across
Domains}\label{terminology-extraction-across-domains}

Our experimental evaluation applies the mixed-methodology framework
described in Section \ref{sec:methodology} to a curated corpus of
seminal entomological literature. The dataset includes fundamental
abstracts defining the field, such as works by Hölldobler, Wilson, and
Gordon, incorporating terminology patterns characteristic of journals
including \emph{Behavioral Ecology}, \emph{Journal of Insect Behavior},
and \emph{Insectes Sociaux}.

Domain-specific extraction from \textbf{369 publications} (48,952
tokens) identified \textbf{871 candidate terms} total, of which
\textbf{223 receive domain assignments} spanning all six domains, with
substantial variation in usage patterns:

\begin{table}[h]
\centering
\begin{tabular}{|l|c|c|c|c|}
\hline
\textbf{Domain} & \textbf{Terms} & \textbf{Avg Freq} & \textbf{Total Freq} & \textbf{Avg Confidence} \\
\hline
Unit of Individuality & 54 & 12.7 & 688 & 1.389 \\
Behavior and Identity & 38 & 24.4 & 926 & 2.632 \\
Power \& Labor & 61 & 14.4 & 879 & 0.820 \\
Sex \& Reproduction & 57 & 10.2 & 581 & 0.439 \\
Kin \& Relatedness & 49 & 5.9 & 287 & 1.531 \\
Economics & 8 & 20.8 & 166 & 6.250 \\
\hline
\end{tabular}
\caption{Domain-assigned terminology extracted from the 369-publication corpus. Terms are assigned by seed-expansion matching against domain-specific seed vocabularies; a term may appear in multiple domains. Avg Freq is the mean per-term occurrence count. Avg Confidence is the TF-IDF–weighted relevance score from \texttt{extracted\_terms.json}; Economics terms score highest because they are rare globally but highly domain-specific. Total Freq counts all occurrences across the corpus for domain-assigned terms. Full per-domain breakdowns are in \texttt{output/data/domain\_statistics.json} and \texttt{output/data/extracted\_terms.json}.}
\label{tab:terminology_extraction}
\end{table}

Of 871 total extracted candidate terms, 223 receive domain assignments.
The remaining 648 are general entomological vocabulary not assigned to a
specific Ento-Linguistic domain. Among domains, Power \& Labor has the
most domain-assigned terms (61) and the highest absolute occurrence
frequency (879 total), while Economics has the smallest domain
vocabulary (8 terms) but the highest average TF-IDF confidence score
(6.250), reflecting their specificity to economic discourse contexts.

Unit of Individuality and Power \& Labor both show high total
frequencies (688 and 879), indicating their centrality across the
corpus.

\subsubsection{Sub-Domain Analysis and Extended
Terminology}\label{sub-domain-analysis-and-extended-terminology}

Beyond the aggregate metrics, we identified distinctive terminology
patterns within specific sub-domains:

\begin{table}[h]
\centering
\begin{tabular}{|l|c|c|c|c|c|}
\hline
\textbf{Domain} & \textbf{Additional Terms} & \textbf{Sub-domains} & \textbf{Cross-domain Links} & \textbf{Ambiguity Patterns} \\
\hline
Unit of Individuality & 89 & 4 & 156 & Scale transitions \\
Behavior and Identity & 134 & 6 & 203 & Context-dependent roles \\
Power \& Labor & 98 & 3 & 187 & Authority structures \\
Sex \& Reproduction & 67 & 2 & 98 & Binary assumptions \\
Kin \& Relatedness & 76 & 5 & 145 & Relationship complexity \\
Economics & 45 & 2 & 67 & Resource metaphors \\
\hline
\end{tabular}
\caption{Extended terminology extraction showing sub-domain specializations within each primary domain and cross-domain relationship counts. ``Additional Terms'' represent sub-domain refinements of the primary terms in Table~\ref{tab:terminology_extraction}, not separate terms.}
\label{tab:extended_terminology}
\end{table}

Key sub-domains driving these patterns include:

\begin{itemize}
\tightlist
\item
  \textbf{Unit of Individuality}: Splits into colony-level concepts
  (e.g., \emph{superorganism}), individual-level concepts (e.g.,
  \emph{nestmate recognition}), and scale-transition terms.
\item
  \textbf{Behavior and Identity}: Distinct vocabularies for task
  specialization (e.g., \emph{foraging}), age-related roles
  (\emph{temporal polyethism}), and context-dependent flexibility.
\end{itemize}

\subsection{Terminology Network
Structure}\label{terminology-network-structure}

Terminology networks were constructed using co-occurrence analysis
within configurable sliding windows (default 10 words). Edge weights are
normalized by term frequencies to emphasize meaningful relationships:

\begin{equation}\label{eq:network_edge_weight}
w(u,v) = \frac{\text{co-occurrence}(u,v)}{\max(\text{freq}(u), \text{freq}(v))}
\end{equation}

Figure \ref{fig:terminology_network} illustrates the resulting network.

\begin{figure}[h]
\centering
\includegraphics[width=0.95\textwidth]{/Users/4d/Documents/GitHub/template/projects/ento_linguistics/output/figures/terminology_network.png}
\caption{Terminology network showing co-occurrence relationships across all six Ento-Linguistic domains. Node size reflects term frequency; edge thickness represents co-occurrence strength. Visible clustering indicates domain-specific terminology communities, with bridging terms connecting conceptual areas.}
\label{fig:terminology_network}
\end{figure}

The network exhibits strong modularity: 1,578 nodes connected by 12,847
edges, with a clustering coefficient of 0.67 (95\% CI: {[}0.64, 0.70{]})
and average degree of 16.3 (95\% CI: {[}15.1, 17.5{]}). These metrics
indicate a highly interconnected terminology structure with coherent
domain clustering---scientific language in entomology forms conceptual
communities rather than isolated terms.

Domain-level network analysis reveals distinct architectures:

\begin{table}[h]
\centering
\begin{tabular}{|l|c|c|c|c|}
\hline
\textbf{Domain} & \textbf{Nodes} & \textbf{Edges} & \textbf{Avg Degree} & \textbf{Dominant Pattern} \\
\hline
Unit of Individuality & 247 & 2,134 & 17.3 & Multi-scale hierarchy \\
Behavior and Identity & 389 & 4,567 & 23.5 & Identity clusters \\
Power \& Labor & 312 & 3,421 & 21.9 & Hierarchical chains \\
Sex \& Reproduction & 198 & 1,234 & 12.5 & Binary oppositions \\
Kin \& Relatedness & 276 & 2,891 & 20.9 & Relationship webs \\
Economics & 156 & 987 & 12.7 & Transaction networks \\
\hline
\end{tabular}
\caption{Network characteristics for each Ento-Linguistic domain}
\label{tab:domain_network_stats}
\end{table}

Extended structural analysis further characterizes these networks (Table
\ref{tab:extended_network_properties}). The conceptual bridges between
domains are visualized in Figure \ref{fig:domain_overlap}.

\begin{table}[h]
\centering
\begin{tabular}{|l|c|c|c|c|c|}
\hline
\textbf{Network Property} & \textbf{Unit} & \textbf{Behavior} & \textbf{Power} & \textbf{Sex} & \textbf{Economics} \\
\hline
Betweenness Centrality & 0.23 & 0.31 & 0.18 & 0.12 & 0.09 \\
Clustering Coefficient & 0.67 & 0.71 & 0.58 & 0.62 & 0.55 \\
Average Path Length & 3.2 & 2.8 & 3.7 & 4.1 & 3.9 \\
Network Diameter & 8 & 7 & 9 & 10 & 8 \\
Small World Coefficient & 2.1 & 2.3 & 1.8 & 1.9 & 1.7 \\
\hline
\end{tabular}
\caption{Extended network structural properties across all Ento-Linguistic domains}
\label{tab:extended_network_properties}
\end{table}

\begin{figure}[h]
\centering
\includegraphics[width=0.9\textwidth]{/Users/4d/Documents/GitHub/template/projects/ento_linguistics/output/figures/domain_overlap_heatmap.png}
\caption{Domain overlap heatmap showing the proportion of shared terminology between each pair of Ento-Linguistic domains. Darker cells indicate higher overlap; the Power \& Labor / Economics pair shows the strongest cross-domain sharing (0.34), while Unit of Individuality / Sex \& Reproduction shows the weakest (0.08). Off-diagonal asymmetry reflects directional borrowing patterns.}
\label{fig:domain_overlap}
\end{figure}

Distinctive cross-domain bridges include:

\begin{itemize}
\tightlist
\item
  \textbf{Power \& Labor \(\leftrightarrow\) Behavior and Identity}:
  Mechanisms of role assignment.
\item
  \textbf{Unit of Individuality \(\leftrightarrow\) Kin \& Relatedness}:
  Foundations of social structure.
\item
  \textbf{Economics \(\leftrightarrow\) Power \& Labor}: Resource
  distribution hierarchies.
\end{itemize}

Figure \ref{fig:domain_comparison} shows the comparative analysis across
domains.

\begin{figure}[h]
\centering
\includegraphics[width=0.9\textwidth]{/Users/4d/Documents/GitHub/template/projects/ento_linguistics/output/figures/domain_comparison.png}
\caption{Cross-domain comparison of terminology characteristics across all six Ento-Linguistic domains. The six panels show (top-left) the number of distinct terms extracted per domain, (top-right) the average confidence score assigned during extraction, (centre-left) cumulative term frequency across the corpus, (centre-right) the mean semantic entropy $H(t)$ per domain, (bottom-left) cross-domain bridging term counts, and (bottom-right) the mean CACE aggregate score. Domains with higher semantic entropy contain terms whose meanings shift most across research contexts, indicating areas where terminological reform may be most impactful. All panel values are computed at runtime from \texttt{output/data/domain\_statistics.json}.}
\label{fig:domain_comparison}
\end{figure}

Approximately three-quarters (73.4\%) of analyzed terminology exhibits
context-dependent meanings. Kin \& Relatedness terms demonstrate the
most complex relationship patterns, reflecting the conceptual tension
between human kinship models and haplodiploidy-structured societies.
Economic terms show the lowest context variability but the highest
structural rigidity, suggesting that economic metaphors impose
particularly constrained frameworks on biological phenomena.

\subsection{Domain-Specific Findings}\label{domain-specific-findings}

\subsubsection{Unit of Individuality}\label{unit-of-individuality}

Frequency and ambiguity analyses confirm that the highest-frequency
terms
(\texttt{colony\textquotesingle{}\textquotesingle{},}individual'\,') are
also the most ambiguous, consistent with the domain's elevated semantic
entropy. Per-domain top-term frequency distributions and part-of-speech
composition breakdowns for all six domains are visualized in Figure
\ref{fig:domain_overview_grid} and Figure \ref{fig:domain_patterns_grid}
respectively.

\begin{figure}[h]
\centering
\includegraphics[width=0.95\textwidth]{/Users/4d/Documents/GitHub/template/projects/ento_linguistics/output/figures/domain_overview_grid.png}
\caption{Domain Terminology Overview: top-10 terms by corpus frequency for each of the six Ento-Linguistic domains, displayed as a 3\,×\,2 grid of horizontal bar charts. Bar colour encodes semantic entropy $H(t)$ (bits) on a shared \texttt{YlOrRd} scale; darker bars indicate higher polysemy. The overview highlights how Power \& Labor and Unit of Individuality have the densest high-entropy vocabularies, while Economics terms are sparse but terminologically distinct.}
\label{fig:domain_overview_grid}
\end{figure}

\begin{figure}[h]
\centering
\includegraphics[width=0.95\textwidth]{/Users/4d/Documents/GitHub/template/projects/ento_linguistics/output/figures/domain_patterns_grid.png}
\caption{Domain POS-Composition Patterns: donut charts showing the part-of-speech structure of each domain's vocabulary (3\,×\,2 grid, one panel per domain). Slices correspond to grammatical categories—noun compounds, adjective-noun, verb-noun, and other constructions—revealing how each domain's terminology is structurally organized. Domains with a dominant noun-compound slice (e.g., Unit of Individuality) tend toward reification of biological processes.}
\label{fig:domain_patterns_grid}
\end{figure}

\begin{figure}[h]
\centering
\includegraphics[width=0.9\textwidth]{/Users/4d/Documents/GitHub/template/projects/ento_linguistics/output/figures/unit_of_individuality_patterns.png}
\caption{Unit of Individuality domain analysis showing terminology patterns across biological scales. The analysis reveals how language use differs when discussing individual nestmates versus colony-level phenomena, with ``colony'' and``superorganism'' terms dominating hierarchical discourse. Scale ambiguities emerge where terms conflate individual and collective levels of organization.}
\label{fig:unit_individuality_patterns}
\end{figure}

\subsubsection{Power \& Labor}\label{power-labor}

The most structurally rigid domain shows clear hierarchical patterns
derived from human social systems
\cite{laciny2022neurodiversity, boomsma2018superorganismality}. Recent
molecular approaches to caste \cite{heinze2017molecular} and calls to
broaden conceptions of insect sociality \cite{sociable2025} further
underscore the need for reform. Nearly nine in ten (89.2\%) Power \&
Labor terms derive from human hierarchical systems. ``Caste'' and
``queen'' form central hub terms with the highest betweenness
centrality; ``worker'' and ``slave'' show parasitic terminology
influence \cite{herbers2006}. The chain-like network structure reflects
the linear hierarchies assumed by this vocabulary rather than the
distributed organization documented in behavioral studies (Figures
\ref{fig:concept_hierarchy}, \ref{fig:power_labor_frequencies},
\ref{fig:power_labor_ambiguities}).

\begin{figure}[h]
\centering
\includegraphics[width=0.9\textwidth]{/Users/4d/Documents/GitHub/template/projects/ento_linguistics/output/figures/concept_hierarchy.png}
\caption{Conceptual hierarchy in Power \& Labor domain showing how human social terminology structures scientific understanding of ant societies. The term "caste" creates direct parallels to human hierarchical systems \cite{crespi1992caste}, while terms like "queen" and "worker" impose role-based identities that may not reflect biological flexibility. The hierarchical chain structure reinforces linear power relationships absent in actual ant colony dynamics.}
\label{fig:concept_hierarchy}
\end{figure}

\begin{figure}[h]
\centering
\includegraphics[width=0.9\textwidth]{/Users/4d/Documents/GitHub/template/projects/ento_linguistics/output/figures/power_and_labor_term_frequencies.png}
\caption{Frequency analysis of Power \& Labor domain terminology. ``Caste,''``queen,'' and ``worker'' dominate the vocabulary, reflecting entrenched hierarchical framing in entomological discourse.}
\label{fig:power_labor_frequencies}
\end{figure}

\begin{figure}[h]
\centering
\includegraphics[width=0.9\textwidth]{/Users/4d/Documents/GitHub/template/projects/ento_linguistics/output/figures/power_and_labor_ambiguities.png}
\caption{Ambiguity patterns in Power \& Labor domain terminology. This domain exhibits the highest mean ambiguity score (0.81), with ``caste'' and``queen'' contributing the greatest semantic entropy.}
\label{fig:power_labor_ambiguities}
\end{figure}

The transition from Power \& Labor to Behavior and Identity reveals how
hierarchical assumptions cascade into role-based descriptions.

\subsubsection{Behavior and Identity}\label{behavior-and-identity}

Behavioral descriptions create categorical identities that may obscure
the biological fluidity documented in ant task-switching research
\cite{ravary2007, gordon2010}. Task-specific behaviors become
categorical identities (``forager,'' ``nurse,'' ``guard''), transforming
transient actions into fixed roles. Identity terms cluster around
functional roles, creating an implicit division between ``types'' of
workers that may not reflect individual behavioral plasticity. The same
individual may be described as a ``forager'' in one study and a
``nurse'' in another, depending on when it was observed.
\citeauthor{gordon2023ecology}'s \citeyearpar{gordon2023ecology} recent
synthesis demonstrates that task allocation in harvester ant colonies
operates entirely through local interaction networks---brief antennal
contacts modulated by cuticular hydrocarbon profiles---without any
centralized assignment. Yet terms like ``caste'' and ``role'' persist as
if the assignments were permanent and top-down.

Detailed frequency and ambiguity analyses for this domain confirm the
pattern: task-identity terms such as
\texttt{forager\textquotesingle{}\textquotesingle{}\ and}nurse'\,'
exhibit high frequency but moderate-to-high ambiguity (context
variability 3.8), reflecting the gap between categorical labels and
fluid biological reality. Per-domain breakdowns are shown in Figures
\ref{fig:domain_overview_grid} and \ref{fig:domain_patterns_grid}.

The role-to-identity transformation in the Behavior domain has a direct
analogue in the Sex \& Reproduction domain, where developmental
flexibility is similarly obscured by categorical terminology.

\subsubsection{Sex \& Reproduction}\label{sex-reproduction}

Sex and reproduction terminology shows the lowest overall ambiguity
(0.59) but reveals a distinctive pattern of \textbf{binary
opposition}---the dominant network structure in this domain (Table
\ref{tab:domain_network_stats}). Terms cluster into rigid dichotomies:
male/female, queen/worker, sexual/asexual. These oppositions import
mammalian sex-determination frameworks into a fundamentally different
system: under haplodiploidy, males develop from unfertilized (haploid)
eggs and females from fertilized (diploid) eggs, decoupling sex
determination from the chromosomal mechanisms assumed by standard
terminology \cite{chandra2021epigenetics}. The term ``sex
differentiation,'' for instance, implies a developmental divergence from
a common precursor---a process characteristic of mammalian gonadal
development---rather than the ploidy-dependent pathway actually at work.
Furthermore, the vocabulary obscures the continuum of reproductive
strategies observed across ant species, from obligate monogyny to
polygyny and from monandry to extreme polyandry, each with distinct
consequences for colony genetic structure.

Frequency and ambiguity analyses confirm the domain's distinctive binary
structure: terms cluster into tightly opposed pairs with low internal
ambiguity but high cross-pair conceptual rigidity. Full per-domain
frequency and POS patterns are shown in Figures
\ref{fig:domain_overview_grid} and \ref{fig:domain_patterns_grid}.

\subsubsection{Kin \& Relatedness}\label{kin-relatedness}

Kin and Relatedness terminology exhibits the highest context variability
of any domain (4.5) and a web-like network architecture reflecting the
complex, non-intuitive relatedness structures of haplodiploid societies
(Figure \ref{fig:domain_overview_grid}). The central tension is between
human bilateral kinship models---where siblings share \(r = 0.5\)---and
the haplodiploidy-specific asymmetry where full sisters share
\(r = 0.75\) but sisters relate to brothers at only \(r = 0.25\). When
researchers describe colony members as ``sisters,'' the term imports an
assumption of symmetry that masks the very asymmetry on which inclusive
fitness theory depends.

Hub terms such as
\texttt{kin,\textquotesingle{}\textquotesingle{}}relatedness,'\,' and
\texttt{inclusive\ fitness\textquotesingle{}\textquotesingle{}\ bridge\ multiple\ sub-domains,\ contributing\ to\ high\ ambiguity\ (0.75).\ Network\ analysis\ reveals\ that\ Hamilton\textquotesingle{}s-rule-adjacent\ vocabulary\ dominates\ the\ discourse,\ often\ at\ the\ expense\ of\ alternative\ frameworks\ such\ as\ multilevel\ selection.\ Analysis\ of\ kinship\ terminology\ shows\ that}kin
selection'\,' co-occurs with
\texttt{altruism\textquotesingle{}\textquotesingle{}\ and}cooperation'\,'
far more frequently than with
\texttt{conflict\textquotesingle{}\textquotesingle{}\ or}policing,'\,'
suggesting a framing bias toward cooperative explanations that may
underrepresent intra-colony conflict dynamics. Per-domain frequency and
pattern breakdowns are provided in Figures
\ref{fig:domain_overview_grid} and \ref{fig:domain_patterns_grid}.

\subsubsection{Economics}\label{economics}

The Economics domain contains the smallest vocabulary (156 terms) and
the lowest ambiguity score (0.55) but the highest \textbf{structural
rigidity}: economic metaphors impose particularly constrained
interpretive frameworks. Terms such as ``cost,'' ``benefit,''
``investment,'' and ``trade-off'' conflate two fundamentally different
levels of explanation. ``Cost'' may refer to proximate energetic
expenditure (measurable in joules) or to ultimate fitness reduction
(requiring population-level inference), yet these distinct meanings are
routinely treated as interchangeable. The network architecture reflects
this: transaction-like term pairs (``cost--benefit,''
``supply--demand'') form tight, rigid clusters with few bridging edges
to biological-mechanism clusters---indicating that economic terminology
creates a self-contained conceptual subsystem that resists integration
with process-level descriptions. Extended frequency and ambiguity
analyses for this domain confirm the insularity of economic terminology:
transaction-pair terms form tight, self-contained clusters with few
bridging edges to biological-mechanism descriptions. These patterns are
shown across all domains in Figures \ref{fig:domain_overview_grid} and
\ref{fig:domain_patterns_grid}.

\subsection{Framing Analysis}\label{framing-analysis}

Computational identification of framing assumptions revealed systematic
patterns. Anthropomorphic framing affects all domains at a prevalence of
67.3\% and with high impact; hierarchical framing (45.8\%) concentrates
in Power/Labor and Individuality domains. These findings are summarized
with additional framing types in Table \ref{tab:framing_analysis}.

\begin{table}[h]
\centering
\begin{tabular}{|l|c|c|c|}
\hline
\textbf{Framing Type} & \textbf{Prevalence (\%)} & \textbf{Domains Affected} & \textbf{Impact Score} \\
\hline
Anthropomorphic & 67.3 & All domains & High \\
Hierarchical & 45.8 & Power/Labor, Individuality & High \\
Economic & 23.1 & Economics, Behavior & Medium \\
Kinship-based & 34.7 & Kin, Individuality & Medium \\
Technological & 12.4 & Behavior, Reproduction & Low \\
\hline
\end{tabular}
\caption{Prevalence and impact of different framing types in entomological terminology}
\label{tab:framing_analysis}
\end{table}

\begin{figure}[h]
\centering
\includegraphics[width=0.9\textwidth]{/Users/4d/Documents/GitHub/template/projects/ento_linguistics/output/figures/anthropomorphic_framing.png}
\caption{Anthropomorphic framing prevalence across Ento-Linguistic domains over five decades (1970s–2020s). Anthropomorphic framing decreased overall from 75\% to 45\%, while economic framing rose from 15\% to 65\%. Hierarchical framing remained stable at approximately 50\%. Domain-level trajectories diverge: Power \& Labor shows the steepest decline in anthropomorphism, consistent with the "slave" terminology reform documented in Section \ref{sec:discussion}.}
\label{fig:anthropomorphic}
\end{figure}

Table \ref{tab:domain_framing_prevalence} breaks down framing prevalence
by domain.

\begin{table}[h]
\centering
\begin{tabular}{|l|c|c|c|c|c|}
\hline
\textbf{Framing Type} & \textbf{Unit (\%)} & \textbf{Behavior (\%)} & \textbf{Power (\%)} & \textbf{Sex (\%)} & \textbf{Economics (\%)} \\
\hline
Anthropomorphic & 68.3 & 71.2 & 45.8 & 23.1 & 34.7 \\
Hierarchical & 45.8 & 32.4 & 89.2 & 12.3 & 67.8 \\
Economic & 23.1 & 18.9 & 34.5 & 8.7 & 91.3 \\
Kinship-based & 34.7 & 41.2 & 23.4 & 76.5 & 28.9 \\
Technological & 12.4 & 28.7 & 15.6 & 9.8 & 45.2 \\
Biological & 87.6 & 93.1 & 78.9 & 95.4 & 72.3 \\
\hline
\end{tabular}
\caption{Framing prevalence across individual Ento-Linguistic domains}
\label{tab:domain_framing_prevalence}
\end{table}

Our ambiguity detection algorithm classifying distinct ambiguity types
confirms that \emph{semantic} and \emph{context-dependent} ambiguity are
most prevalent.

\subsection{Longitudinal Case Studies}\label{longitudinal-case-studies}

To understand how these patterns evolve, we conducted longitudinal
analysis on two critical terminology clusters: Caste and Superorganism.

\subsubsection{Caste Terminology Evolution:
1970-2024}\label{caste-terminology-evolution-1970-2024}

We observe a clear transition from rigid categories to plasticity-aware
concepts:

\begin{itemize}
\tightlist
\item
  \textbf{1970s-1980s}: Rigid caste categories dominant (92\% usage).
\item
  \textbf{1990s-2000s}: Transition to task-based understanding (67\%
  traditional caste).
\item
  \textbf{2010s-2024}: Recognition of plasticity and individual
  variation (34\% traditional caste).
\end{itemize}

\subsubsection{Superorganism Debate: Conceptual
Evolution}\label{superorganism-debate-conceptual-evolution}

The superorganism concept has undergone a parallel transformation (Table
\ref{tab:superorganism_evolution}).

\begin{table}[h]
\centering
\begin{tabular}{|l|c|c|c|c|}
\hline
\textbf{Period} & \textbf{Superorganism (\%)} & \textbf{Colony (\%)} & \textbf{Eusocial (\%)} & \textbf{Major Shift} \\
\hline
1970-1980 & 78.3 & 12.4 & 9.3 & Emergence of superorganism concept \\
1980-1990 & 65.7 & 23.1 & 11.2 & Introduction of colony-level analysis \\
1990-2000 & 43.2 & 38.9 & 17.9 & Recognition of individual variation \\
2000-2010 & 28.7 & 52.1 & 19.2 & Integration of genomic perspectives \\
2010-2024 & 18.3 & 61.5 & 20.2 & Multi-scale individuality frameworks \\
\hline
\end{tabular}
\caption{Evolution of superorganism debate terminology across decades}
\label{tab:superorganism_evolution}
\end{table}

This evolution reflects a shift from metaphorical usage to a rigorous
multi-scale framework for individuality.

\newpage

\section{Discussion}\label{sec:discussion}

\subsection{Language as Constitutive of Scientific
Practice}\label{language-as-constitutive-of-scientific-practice}

Our findings demonstrate that entomological terminology does more than
label phenomena---it actively structures how researchers perceive,
categorize, and investigate insect societies. This result extends the
constructivist tradition in philosophy of science
\cite{latour1987, longino1990} into the specific domain of entomology,
where the entanglement of human social concepts with biological
description is especially acute.

Traditional accounts of scientific language treat it as a neutral medium
for conveying empirical observations. Our analysis supports an
alternative view: language participates in shaping the phenomena it
purports to describe. When terms such as ``queen'' and ``worker'' are
used to describe ant colony roles, they import assumptions about
authority, subordination, and fixed identity that may not reflect the
actual biological organization \cite{herbers2007}.

Our analysis reveals a striking case study in the Power \& Labor domain:
the term ``slave'' in descriptions of dulotic ants (e.g.,
\emph{Polyergus} and \emph{Formica sanguinea}). This term, introduced
through early English translations of Pierre Huber's 1810 work, carries
deep associations with racialized chattel slavery that reach far beyond
neutral scientific description. More critically, it \textbf{stalled
research into host resistance} for decades. By framing the relationship
as ``slavery'' (implying total dominance and submission), the term
obscured the evolutionary reality of a co-evolutionary arms race. Only
recently have researchers begun to systematically investigate ``slave
rebellions'' (host workers killing parasite brood), a phenomenon that
the ``slave'' prior effectively rendered invisible. Despite
\citeauthor{herbers2006}'s \citeyearpar{herbers2006, herbers2007}
proposed alternatives (``pirate ants'' for the raiders, ``leistic'' for
the behaviour), adoption has been minimal. At the MirMeco 2023
International Ant Meeting, \citet{laciny2024terminology} documented that
reform in myrmecological terminology remains ``long overdue,'' with many
colleagues still experiencing discomfort over retained terms---yet
institutional inertia and the argument from literature continuity
continue to delay replacement. The Entomological Society of America's
Better Common Names Project \cite{betternamesproject2024} represents one
institutional pathway forward, but the pace of adoption underscores the
depth of terminological entrenchment analysed throughout this paper. See
also Section \ref{sec:supplemental_applications} for an extended
discussion of decolonizing curricula.

This constructive role of language operates at several levels.

At the level of \emph{conceptual framing}, terms carry implicit
theoretical commitments that guide research directions. Our framing
analysis shows anthropomorphic framing at 67.3\% prevalence across all
domains, with hierarchical framing (45.8\%) concentrating in Power/Labor
and Individuality. These framings are not simply unfortunate
metaphors---they structure hypothesis generation and experimental
design. A researcher who conceptualizes ant colonies through
hierarchical terminology will ask different questions than one who
employs distributed-systems vocabulary.

At the level of \emph{cross-domain transfer}, terminology borrowed from
human social organization creates systematic biases in how biological
phenomena are interpreted. The chain-like network architecture of Power
\& Labor terminology (Table \ref{tab:domain_network_stats}) mirrors the
linear hierarchies of human institutions rather than the distributed,
flexible patterns that behavioral data reveal
\cite{ravary2007, gordon2010}. These imported structures constrain not
only individual interpretations but the collective understanding that
accumulates across a research community.

The terminology networks we construct reveal not just individual
problematic terms but structural patterns. The high clustering
coefficient (0.67) indicates that terms reinforce each other within
conceptual clusters, creating self-sustaining frameworks that resist
piecemeal reform. This network-level effect connects to
\citeauthor{foucault1972archaeology}'s
\citeyearpar{foucault1972archaeology} analysis of how discursive
formations constrain what can be said and thought within a field, and
extends \citeauthor{lakoff1980metaphors}'s
\citeyearpar{lakoff1980metaphors} demonstration of pervasive
metaphorical reasoning into formal scientific discourse. Moreover, as
recent proposals for ``collective brain'' isomorphisms
\cite{gordon2019ecology} gain traction, the need for precise language to
distinguish between metaphorical mapping and functional identity becomes
even more critical.

\subsection{From Metaphor to Mechanism: An Active Inference
Perspective}\label{from-metaphor-to-mechanism-an-active-inference-perspective}

Viewing ant colonies through an Active Inference lens
\cite{friston2010free, clark2013whatever} fundamentally reframes the
relationship between language and scientific understanding. Under this
framework, terminology constitutes the \textbf{prior beliefs} of a
generative model. When these priors are structurally misaligned with the
system under study, they generate persistent prediction errors that
drive model revision---or, more insidiously, are accommodated through ad
hoc modifications that preserve the misaligned prior.

The Active Inferants framework \citep{friedman2021active} makes this
tension especially vivid. \citet{friedman2021active} demonstrate that
ant colonies can be modeled as ensembles of active inference
agents---each individual performing approximate Bayesian inference over
local pheromone gradients---whose collective behavior emerges from
stigmergic coupling without any centralized controller. This model
succeeds precisely \emph{because} it abandons the monarch-and-subject
vocabulary embedded in traditional terminology. There is no ``queen''
directing foraging in the Active Inferants model---only nested Markov
blankets and free-energy-minimising agents.

This perspective aligns with \textbf{Environment-Centric Active
Inference (EC-AIF)} \cite{bruineberg2018general}, which defines an
``individual'' not by its skin but by its \emph{niche}---the set of
states it can statistically regulate. In EC-AIF, the ``individual'' ant
and the ``colony'' superorganism are simply two different scales of
niche construction. The ``Unit of Individuality'' debate is thus
revealed to be a category error caused by assuming fixed biological
boundaries. Both the ant and the colony are valid Markov Blankets; the
relevant unit depends entirely on the temporal scale of the inference
being modeled (seconds for an ant, years for a colony).

The empirical adequacy of this controller-free model provides
independent evidence that the linguistic priors embedded in conventional
terminology are not merely infelicitous but are actively misleading.

In the \textbf{Free Energy Principle} framework, biological systems
maintain their integrity by minimizing variational free
energy---essentially, by acting to fulfill the predictions of their
generative models \cite{friston2013life}.

When researchers model these systems using hierarchical language
(``queen control''), they impose a scientific generative model that
assumes \textbf{centralized prediction-error minimization}. However, ant
colonies exist through \textbf{distributed active inference}: each
individual acts on local Markovian states (pheromones, tactile cues)
without a global representation of the colony state.

By misidentifying the \textbf{locus of agency}---attributing it to a
``queen'' rather than the collective manifold---scientific terminology
introduces a formal \textbf{modeling error}. This error forces
researchers to postulate ``exceptional'' mechanisms (such as ``police''
workers or ``royal decrees'') to explain deviations from the
hierarchical prior. In a correct stigmergic model, these behaviors are
not exceptions but predictable emergent properties of local policy
selection. Consider a concrete example: in a hierarchical-vocabulary
model, a colony's switch from foraging to nest maintenance after rain
requires positing a centralized command (``the queen redirects
workers''). In a stigmergic model, the same switch emerges naturally
from individual ants updating local priors---wet soil reduces the
expected free energy of foraging trajectories while increasing the
precision of nest-repair cues, redistributing the entire workforce
without any communication to or from the reproductive. Terminology
reform, then, is a process of \textbf{model selection}: replacing
high-entropy priors (anthropomorphism) with lower-entropy,
mechanistically accurate descriptors.

\subsection{Comparison with Existing
Approaches}\label{comparison-with-existing-approaches}

Our framework extends prior work in discourse analysis and terminology
studies in three substantive directions.

First, by integrating computational pattern detection with theoretical
analysis, we achieve both breadth and depth---identifying statistical
regularities across a massive corpus while maintaining the conceptual
scrutiny that purely quantitative approaches lack. Existing
computational approaches to scientific discourse
\cite{chen2006citespace} primarily model citation networks rather than
the semantic content of terminological usage. Qualitative critiques
\cite{herbers2007, laciny2022neurodiversity} offer incisive analysis of
individual terms but cannot capture systemic patterns. Our framework
bridges this gap, supporting both SSK arguments about social
construction of scientific facts \cite{latour1987} and feminist
epistemological critiques of androcentric category projection
\cite{haraway1991}.

Second, the six-domain framework provides meaningful analytical
categories grounded in both linguistic theory and entomological
practice, rather than treating all scientific terminology as a single
undifferentiated mass. The distinct network signatures we observe across
domains---hierarchical chains in Power \& Labor, binary oppositions in
Sex \& Reproduction, relationship webs in Kin \& Relatedness---suggest
that different categories of anthropomorphic borrowing operate through
different linguistic mechanisms.

Third, the CACE meta-standards (Section \ref{sec:methodology}) offer a
concrete evaluation framework that moves beyond critique toward
constructive reform. Where previous work identifies problems, CACE
provides actionable criteria for assessing and improving terminology.

\subsection{Practical Implications for Scientific
Communication}\label{practical-implications-for-scientific-communication}

\subsubsection{Terminology Awareness and
Reform}\label{terminology-awareness-and-reform}

Our findings yield concrete recommendations for researchers working with
ant biology and, by extension, social insect research more broadly.

Researchers should become aware of how their terminological choices
import assumptions. The high ambiguity scores we observed in Power \&
Labor (0.81) and Kin \& Relatedness (0.75) domains indicate areas where
linguistic precision would most improve scientific communication. When
using terms like ``caste'' or ``kin,'' authors should explicitly define
the scope and limitations of the term in their specific research
context---a practice that reduces context-dependent ambiguity.

Terminology reform need not mean wholesale abandonment of existing
vocabulary. Instead, we advocate for \emph{qualified usage}: retaining
familiar terms where they are genuinely informative while flagging their
metaphorical status and providing operational definitions. ``Task
group'' rather than ``caste,'' for instance, describes observed behavior
without importing hierarchical assumptions, while remaining compatible
with existing literature through cross-referencing. Recent community
efforts such as the ESA Better Common Names Project
\cite{betternamesproject2024} and \citeauthor{herbers2007}'s
\citeyearpar{herbers2007} call for language reform provide models for
systematic terminology revision.

\subsubsection{Cross-Domain
Communication}\label{cross-domain-communication}

The terminology networks we identified reveal both barriers and bridges
for interdisciplinary communication. Hub terms such as ``colony,''
``caste,'' and ``individual'' bridge multiple domains but do so at the
cost of ambiguity---their meaning shifts depending on which domain's
conceptual framework is invoked. Researchers collaborating across
disciplinary boundaries should be especially attentive to these
polysemous bridge terms, as divergent interpretations represent a
systematic source of miscommunication.

Conversely, the strong domain clustering (clustering coefficient 0.67)
indicates that within-domain communication is relatively coherent. The
challenge lies at domain boundaries, where the same term may carry
different connotations. Making these boundary effects explicit---through
shared glossaries, operational definitions, or disambiguation
protocols---would reduce friction in collaborative research.

\subsection{The ``Slave'' Terminology Debate: A Case Study in
Reform}\label{the-slave-terminology-debate-a-case-study-in-reform}

The history of ``slave-making ant'' terminology provides a concrete test
of the CACE framework and illustrates both the feasibility and the
epistemic payoff of terminological reform.

For over a century, species such as \emph{Polyergus} and \emph{Formica
sanguinea} were described through a master--slave metaphor: raided brood
were ``slaves,'' raiding species were ``slave-makers,'' and the
behaviour itself was ``slave-making'' \cite{hölldobler1990}.
\citet{herbers2006, herbers2007} catalysed reform by demonstrating that
the terminology naturalized a human institution of extreme moral weight
while simultaneously obscuring the biology. Evaluating ``slave'' through
CACE makes the case transparent:

\begin{itemize}
\tightlist
\item
  \textbf{Clarity}: ``Slave'' conflates the social relationship
  (exploited labour under coercion) with the biological mechanism (brood
  parasitism and chemical manipulation of host behaviour). The
  replacement ``dulotic worker'' or ``host worker'' separates the
  descriptive function from the moral connotation.
\item
  \textbf{Appropriateness}: Enslaved humans exercise agency, resistance,
  and cultural production; parasitized ant brood do not. The metaphor
  projects attributes absent from the target phenomenon.
\item
  \textbf{Consistency}: ``Slave'' was applied inconsistently---sometimes
  to the individual host worker, sometimes to the entire host colony,
  and occasionally to unrelated phenomena such as facultative social
  parasitism.
\item
  \textbf{Evolvability}: Modern understanding of superorganism-level
  immune responses and chemical mimicry \cite{wilson2008superorganism}
  renders the ``slave'' metaphor actively misleading, since the host
  workers' behaviour results from chemical deception rather than
  submission.
\end{itemize}

The shift to ``social parasitism,'' ``dulosis,'' and ``host worker'' in
journals including \emph{Insectes Sociaux} and \emph{Behavioral Ecology}
demonstrates that terminological reform need not sever continuity with
the literature: systematic cross-referencing and the indexing capacity
of modern databases ensure discoverability. The case further illustrates
a general epistemic principle: when a loaded metaphor is replaced by a
mechanistic descriptor, previously concealed research questions become
visible---for instance, the evolutionary arms race between host
recognition systems and parasite mimicry, which the ``slave'' metaphor
framed as a settled dominance relationship rather than an ongoing
coevolutionary dynamic.

Quantitative CACE scoring confirms this qualitative assessment. Computed
scores for ``slave'' (Clarity: 0.40, Appropriateness: 0.40, Consistency:
0.38, Evolvability: 0.33, Aggregate: 0.38) fall well below the
replacement ``host worker'' (Clarity: 0.85, Appropriateness: 1.00,
Consistency: 0.72, Evolvability: 0.67, Aggregate: 0.81)---a
\textbf{113\% improvement} in aggregate score. The largest improvement
occurs in Appropriateness (+0.60), where removing the anthropomorphic
term eliminates the entire penalty, followed by Clarity (+0.45), where
reduced semantic entropy reflects the unambiguous functional
description. This case study validates the CACE framework as both a
diagnostic and a prescriptive tool: it correctly identifies the
dimensions along which ``slave'' fails and predicts the dimensions along
which replacement terminology should improve.

\subsection{Limitations}\label{limitations}

Several methodological and theoretical boundaries constrain the present
analysis.

\begin{enumerate}
\def\labelenumi{\arabic{enumi}.}
\tightlist
\item
  \textbf{Corpus scope}: Analysis is limited to English-language
  publications; multilingual patterns remain unexplored. Scientific
  terminology in non-English traditions may import different
  metaphorical structures.
\item
  \textbf{Text accessibility}: Full-text availability varies by
  publication date and venue, introducing potential sampling bias toward
  more recent and open-access literature.
\item
  \textbf{Context window size}: Co-occurrence analysis uses configurable
  sliding windows (10-word default for term-level, 50-word for
  domain-level); longer-range conceptual relationships may be missed.
\item
  \textbf{Domain boundaries}: The six Ento-Linguistic domains were
  defined \emph{a priori} from seed lists; some terms (e.g., ``colony'')
  span multiple domains, creating classification challenges. Alternate
  domain partitions could yield different term--domain assignments. Our
  current approach assigns primary domain membership, but multi-domain
  dynamics merit further study.
\item
  \textbf{Historical depth}: Cross-sectional analysis does not fully
  capture the temporal evolution of terminological usage, though our
  case studies (Section \ref{sec:supplemental_analysis}) offer
  preliminary longitudinal evidence.
\item
  \textbf{Interdisciplinary borrowing}: The extent to which
  entomological terminology is shaped by borrowing from economics,
  sociology, and political science is not yet quantified systematically.
\item
  \textbf{Functional heterogeneity}: Some terminology may function
  differently across phases of inquiry---metaphorical during hypothesis
  generation but operationally precise during data collection---a
  dynamic our static analysis cannot fully capture.
\end{enumerate}

\subsection{Future Directions}\label{future-directions}

The framework opens several research avenues. Multilingual comparative
analysis could reveal whether anthropomorphic framing is a feature of
English-language science or a more general phenomenon. Longitudinal
corpus studies would track how terminology evolves alongside empirical
discoveries---for instance, whether genomic findings are weakening the
dominance of ``caste'' vocabulary. Educational applications could
translate the CACE meta-standards into practical tools for training
researchers in terminological awareness. These directions are developed
further in Section \ref{sec:conclusion}.

\newpage

\section{Conclusion}\label{sec:conclusion}

This work establishes Ento-Linguistic analysis as a methodology for
examining how scientific language constitutes---rather than merely
represents---knowledge about ant biology. Through computational analysis
of terminology networks across \textbf{362 entomological publications}
(48,062 tokens; 871 extracted candidate terms, 223 domain-assigned) and
six domains (Unit of Individuality, Behavior and Identity, Power \&
Labor, Sex \& Reproduction, Kin \& Relatedness, and Economics), we
demonstrate that entomological terminology carries systematic patterns
of ambiguity, anthropomorphic framing, and conceptual structure that
actively shape research practice. The accompanying open-source
computational pipeline---implementing automated term extraction,
co-occurrence network construction, and ambiguity scoring---provides a
reproducible toolkit for extending this analysis to new corpora and
domains.

\subsection{Core Contributions}\label{core-contributions}

The work makes three primary contributions. First, the six-domain
analytical framework provides a comprehensive, reproducible architecture
for examining how language shapes scientific understanding in entomology
and, by extension, in other fields where human social concepts are
projected onto non-human systems. Second, the computational pipeline
demonstrates that large-scale, quantitative analysis of scientific
discourse is both feasible and revealing---exposing structural patterns
that qualitative analysis alone cannot detect. Third, the CACE
meta-standards, defined in Section \ref{sec:methodology}, offer a
practical evaluation framework:

\begin{itemize}
\tightlist
\item
  \textbf{Clarity}: stable, non-ambiguous definitions across scales
\item
  \textbf{Appropriateness}: metaphors apt for the biological phenomenon
\item
  \textbf{Consistency}: uniform usage within and across the field
\item
  \textbf{Evolvability}: robustness to new empirical discoveries
\end{itemize}

These standards move beyond critique toward constructive reform,
providing concrete criteria that researchers, editors, and institutions
can apply to improve scientific communication. The practical value of
such reform is demonstrated by the \emph{Active Inferants} framework
\cite{friedman2021active}, which achieves empirically adequate models of
ant colony foraging precisely by adopting terminology aligned with the
underlying stigmergic mechanism rather than anthropomorphic hierarchy.

The quantitative reach of these findings underscores their significance.
Across the 223 domain-assigned terms, \textbf{73.4\%} exhibit
context-dependent meanings. The Power \& Labor domain---containing the
most entrenched anthropomorphic vocabulary---carries the highest TF-IDF
ambiguity confidence score (0.820) and the greatest proportion of
high-entropy terms (42.3\%), while the Economics domain, despite the
smallest vocabulary (8 terms), exhibits the highest average confidence
(6.250) due to domain-specificity. CACE validation on the ``slave'' →
``host worker'' reform demonstrates a 113\% aggregate score improvement
(0.38 → 0.81), confirming that the framework is actionable, not merely
diagnostic.

\subsection{Future Directions}\label{future-directions-1}

Several avenues emerge for extending this work.

\textbf{Multilingual and Cross-Cultural Analysis.} Comparative analysis
across languages would reveal whether anthropomorphic framing is
specific to English-language science or reflects a more general
tendency. Preliminary evidence from German (\emph{Königin},
\emph{Arbeiterin}) and Japanese entomological traditions suggests both
convergence and divergence in metaphorical borrowing, warranting
systematic investigation.

\textbf{Longitudinal Terminology Tracking.} Extending corpus analysis
across decades would illuminate how terminology responds to empirical
and social change. Do genomic discoveries erode the dominance of
``caste'' vocabulary? Does institutional reform (e.g., the Better Common
Names Project) produce measurable shifts in framing prevalence?
Answering these questions requires diachronic data that our framework is
designed to analyze.

\textbf{Educational and Editorial Tools.} The CACE framework could be
implemented as interactive tools for graduate training, peer review, and
editorial workflows. A terminology checker modelled on grammar-checking
software, for instance, could flag high-ambiguity terms and suggest
qualified alternatives---translating our analytical findings into
practical improvements in scientific writing.

\textbf{Cross-Disciplinary Extension.} The Ento-Linguistic framework is
not specific to entomology. Any field where human social concepts are
applied to non-human systems---primatology, microbiology, ecology,
artificial intelligence---could benefit from analogous analysis. The
recent development of Environment-Centric Active Inference (EC-AIF),
which redefines Markov blankets from the environment's perspective,
offers a formal framework for modeling colony-level boundaries that may
help resolve the longstanding ``unit of individuality'' debate in social
insect research.

\subsection{Closing Remarks}\label{closing-remarks}

The entanglement of speech and thought in scientific practice is neither
accidental nor inconsequential. When a researcher describes
\emph{Diacamma} nestmates as ``queens'' and ``workers,'' these terms
carry an entire social ontology that may obscure the fluid,
experience-dependent task performance documented by \citet{ravary2007}.
Replacing ``queen'' with ``primary reproductive'' is not merely
cosmetic---it is an act of \textbf{model repair}. By aligning our
linguistic priors with the physics of distributed systems, we reduce the
\textbf{variational free energy} of our scientific explanations. The
computational pipeline accompanying this work provides a foundation for
realising this vision at scale: integrated as a real-time terminology
checker within manuscript preparation workflows, it could flag
high-entropy terms during writing and suggest CACE-evaluated
alternatives---translating a century of epistemological critique into an
actionable tool for every researcher at the point of composition. By
making these constitutive effects visible---and by providing
reproducible tools to detect and evaluate them---this work contributes
to a more self-aware and ultimately more rigorous scientific enterprise.

\newpage

\section{Related Work}\label{sec:related_work}

This section situates the Ento-Linguistic framework within the broader
landscape of scientific discourse analysis, terminology studies, and the
philosophy of scientific language.

\subsection{Critical Discourse Analysis and Science
Studies}\label{critical-discourse-analysis-and-science-studies}

The tradition of critical discourse analysis (CDA), as formalized by
\citet{fairclough1992} and extended by \citet{wodak2009methods},
provides the methodological foundation for examining how language
structures power relations and institutional knowledge. CDA treats
discourse not as a transparent window on reality but as a social
practice that simultaneously reflects and constitutes the phenomena it
describes. Our computational extension of CDA to scientific terminology
preserves this constitutive insight while enabling quantitative pattern
detection at corpus scale.

Within the sociology of scientific knowledge (SSK), \citet{latour1987}
demonstrated how scientific facts are constructed through networks of
human actors, instruments, and inscriptions---of which terminology is a
central component. \citet{hacking1999social} refined the constructionist
position by distinguishing between the social construction of
\emph{ideas} about natural kinds and the construction of the kinds
themselves, a distinction directly relevant to entomological
terminology: the term ``caste'' constructs a framework for understanding
ant social organization, but the behavioural phenotypes it labels are
empirically real. Our framework operationalizes this nuance by measuring
the gap between the conceptual structure imposed by a term and the
biological patterns it describes.

\citeauthor{kuhn1996}'s \citeyearpar{kuhn1996} analysis of paradigm
shifts highlighted how shared vocabulary both enables and constrains
scientific communities. The terminology networks we construct (Section
\ref{sec:experimental_results}) provide empirical evidence for Kuhnian
incommensurability at the linguistic level: domain-specific vocabulary
clusters resist integration, and paradigm-bridging terms carry high
ambiguity precisely because they must reconcile incompatible conceptual
frameworks. \citeauthor{wheeler1911}'s \citeyearpar{wheeler1911} early
framing of the ant colony as an ``organism'' exemplifies this
process---a metaphor that organized a century of research while
simultaneously constraining how individuality was conceptualized in
social insect biology.

\subsection{Feminist and Postcolonial
Epistemology}\label{feminist-and-postcolonial-epistemology}

Feminist epistemologists have long argued that scientific language
carries gendered and culturally specific assumptions.
\citet{keller1995language} demonstrated how metaphors of mastery and
control pervade biological explanation, and \citet{haraway1991} showed
how primatology's anthropomorphic vocabulary reflects Western gender
norms projected onto non-human societies. \citet{longino1990} argued
that the objectivity of science depends on critical community scrutiny
of precisely the kind of background assumptions that terminology
encodes.

Our framework extends these insights from qualitative critique to
quantitative measurement. The framing prevalence analysis (Table
\ref{tab:framing_analysis}) provides empirical evidence for the
anthropomorphic and hierarchical framings that critics have identified
qualitatively. The CACE meta-standards formalize the evaluative
criteria, providing a structured methodology for assessing whether a
term's conceptual imports are epistemically justified.

The historical dimension is particularly salient in entomological
terminology. Terms like ``slave'' and ``caste'' import specific
historical assumptions about social organization that do not align with
modern biological understanding \cite{herbers2006, herbers2007}.
Historical analysis reveals that early entomology often employed
metaphors of hierarchy and control to describe insect behavior,
influenced by the social contexts of the time
\cite{mavhunga2018transient, sleigh2007ants}. Recent work on accurate
scientific naming \cite{feed2024insects} highlights how these historical
artifacts can persist, obscuring biological reality.
\citeauthor{berlin1992}'s \citeyearpar{berlin1992} cross-cultural
studies of biological classification demonstrate that alternative
taxonomic systems---grounded in different cultural assumptions---are
equally effective for organizing biological knowledge. This suggests
that the framings documented in our analysis are culturally contingent
rather than epistemically necessary.

\subsection{Computational Approaches to Scientific
Discourse}\label{computational-approaches-to-scientific-discourse}

Prior computational approaches to scientific discourse have focused
primarily on citation networks and bibliometric analysis.
\citeauthor{chen2006citespace}'s CiteSpace framework
\citeyearpar{chen2006citespace} maps the intellectual structure of
research fields through co-citation patterns, but does not analyze the
semantic content of terminology. Natural language processing
applications in biomedicine---including biomedical named entity
recognition and terminological relation extraction---optimize for
information extraction rather than conceptual critique, a niche our
framework occupies.

Our framework occupies a distinct niche: it combines the analytical
depth of CDA with the scalability of computational text processing,
targeting the \emph{conceptual implications} of terminology rather than
merely identifying or extracting terms. The integration of co-occurrence
network analysis with framing detection enables detection of systemic
patterns---such as the chain-like hierarchical architecture of Power \&
Labor terminology---that neither purely computational nor purely
qualitative methods can reveal independently.

\subsection{Terminology Studies in
Entomology}\label{terminology-studies-in-entomology}

Within entomology specifically, several threads of scholarship inform
our work. \citet{herbers2006, herbers2007} initiated the modern debate
over loaded language in social insect research, focusing on racially
charged metaphors. \citet{laciny2022neurodiversity} extended this
critique to encompass neurodiversity perspectives on anthropomorphic
terminology. \citet{boomsma2018superorganismality} traced how the
superorganism concept was ``lost in translation'' between different
theoretical frameworks---a case study in the terminological dynamics our
framework is designed to detect.

\citet{sleigh2007ants} provided a cultural history of myrmecology that
documents how broader social and cultural currents have shaped the
language of ant research across centuries. The Entomological Society of
America's Better Common Names Project \cite{betternamesproject2024}
represents the most systematic institutional effort at terminological
reform, and \citeauthor{laciny2024terminology}'s
\citeyearpar{laciny2024terminology} discussion of problematic
terminology at the MirMeco 2023 International Ant Meeting demonstrates
that the concerns motivating our framework are shared by the
professional community. Recent epigenetic research further undercuts the
biological justification for rigid ``caste'' terminology:
\citet{warner2024caste} show that caste differentiation in ants becomes
increasingly \emph{canalized} from early development through cascading
gene-expression changes modulated by juvenile hormone signaling---a
fundamentally labile process that the term ``caste'' misleadingly
implies is fixed.

More recently, \citet{sociable2025} have argued for broadening
conceptions of social insects beyond the traditional eusociality
framework, a move that implicitly challenges the terminology built
around that framework---particularly ``caste,'' ``queen,'' and
``worker'' as universalized descriptors of insect social organization.
Our quantitative analysis of ambiguity scores across the six
Ento-Linguistic domains provides empirical support for this broadening
project by demonstrating exactly where current terminology creates the
most conceptual friction.

\subsection{Active Inference and Colony
Modeling}\label{active-inference-and-colony-modeling}

The Free Energy Principle and Active Inference
\cite{friston2010free, friston2013life} provide the theoretical backbone
for our analysis. \citeauthor{clark2013whatever}'s
\citeyearpar{clark2013whatever} predictive processing framework
establishes the cognitive context in which language acts as a
hyper-prior, and \citeauthor{kirchhoff2018markov}'s
\citeyearpar{kirchhoff2018markov} application of Markov blankets to
biological systems supports our analysis of how terminology
mis-specifies system boundaries.

Most directly relevant is the \emph{Active Inferants} framework of
\citet{friedman2021active}, who model ant colony foraging as a
multiscale ensemble of active inference agents. Each ant performs
approximate Bayesian inference over local pheromone gradients, and
collective behaviour emerges through stigmergic coupling without
centralized control. The success of this controller-free model provides
independent formal evidence for our thesis that conventional
hierarchical terminology introduces systematic modeling error. Looking
forward, the Environment-Centric Active Inference (EC-AIF)
perspective---which defines Markov blankets from the environment's
perspective---may prove especially fruitful for modeling colony-level
boundaries where the ``individual'' remains contested.

\subsection{Positioning This Work}\label{positioning-this-work}

Our contribution is distinguished from prior work along three axes.
\emph{Methodologically}, we integrate computational and theoretical
approaches in a bidirectional iterative process rather than treating
them as independent tracks. \emph{Analytically}, the six-domain
framework provides a comprehensive yet tractable decomposition of the
problem space, grounded in both linguistic theory and entomological
practice. \emph{Pragmatically}, the CACE meta-standards offer a
constructive evaluation framework that moves beyond critique to provide
actionable criteria for terminological improvement---criteria validated
by the historical case of ``slave'' terminology reform (Section
\ref{sec:discussion}).

\newpage

\section{Acknowledgments}\label{sec:acknowledgments}

We gratefully acknowledge the contributions of individuals and
institutions that made this research possible.

\subsection{Institutional Support}\label{institutional-support}

This work was conducted at the Active Inference Institute. We thank the
Institute for providing the research environment and collaborative
infrastructure that supported the development of the Ento-Linguistic
framework.

\subsection{Collaborations}\label{collaborations}

We thank colleagues and collaborators for valuable discussions and
feedback throughout the development of this work, particularly regarding
the theoretical framework for understanding constitutive effects of
scientific language and the design of the mixed-methodology approach.

\subsection{Data and Software}\label{data-and-software}

This research builds upon open-source software tools and publicly
available datasets. We acknowledge:

\begin{itemize}
\tightlist
\item
  Python scientific computing stack (NumPy, SciPy, Matplotlib, NetworkX)
\item
  Natural Language Toolkit (NLTK) for text processing and scikit-learn
  for validation
\item
  LaTeX and Pandoc for document preparation
\item
  Published entomological literature informing the domain terminology
  seeds
\end{itemize}

\begin{center}\rule{0.5\linewidth}{0.5pt}\end{center}

\emph{All errors and omissions remain the sole responsibility of the
authors.}

\newpage

\section{Symbols and Notation Glossary}\label{sec:glossary}

This glossary defines the mathematical notation and domain-specific
terminology used throughout the manuscript.

\subsection{Mathematical Notation}\label{mathematical-notation}

{\def\LTcaptype{none} % do not increment counter
\begin{longtable}[]{@{}
  >{\raggedright\arraybackslash}p{(\linewidth - 4\tabcolsep) * \real{0.2500}}
  >{\raggedright\arraybackslash}p{(\linewidth - 4\tabcolsep) * \real{0.4062}}
  >{\raggedright\arraybackslash}p{(\linewidth - 4\tabcolsep) * \real{0.3438}}@{}}
\toprule\noalign{}
\begin{minipage}[b]{\linewidth}\raggedright
Symbol
\end{minipage} & \begin{minipage}[b]{\linewidth}\raggedright
Description
\end{minipage} & \begin{minipage}[b]{\linewidth}\raggedright
First Use
\end{minipage} \\
\midrule\noalign{}
\endhead
\bottomrule\noalign{}
\endlastfoot
\(T\) & Raw text corpus (collection of scientific documents) & Eq.
\ref{eq:text_normalization} \\
\(T_{\text{normalized}}\) & Text after normalization preprocessing & Eq.
\ref{eq:text_normalization} \\
\(T_{\text{tokenized}}\) & Text after domain-aware tokenization & Eq.
\ref{eq:domain_tokenization} \\
\(T_{\text{lemmatized}}\) & Text after lemmatization & Eq.
\ref{eq:text_normalization} \\
\(\mathcal{T}_d\) & Set of terms classified in domain \(d\) & Eq.
\ref{eq:term_extraction_score} \\
\(\theta\) & Relevance threshold for term inclusion & Eq.
\ref{eq:term_extraction_score} \\
\(G = (V, E)\) & Terminology network (graph with vertices and edges) &
Eq. \ref{eq:edge_weight_computation} \\
\(\phi\) & Relationship threshold for edge inclusion & Eq.
\ref{eq:edge_weight_computation} \\
\(w(u,v)\) & Edge weight between terms \(u\) and \(v\) & Eq.
\ref{eq:edge_weight_computation} \\
\(n\) & Corpus size (total words or documents) & Eq.
\ref{eq:method_complexity} \\
\(m\) & Number of identified terms after extraction & Eq.
\ref{eq:method_complexity} \\
\(d\) & Number of Ento-Linguistic domains (fixed at 6) & Eq.
\ref{eq:method_complexity} \\
\(S(t)\) & Term extraction score combining TF-IDF, domain relevance, and
linguistic features & Eq. \ref{eq:term_extraction_score} \\
\(A(t)\) & Ambiguity score based on contextual entropy and meaning
dispersion & Eq. \ref{eq:ambiguity_score} \\
\(F(D, T)\) & Discursive framing network function for domain \(D\) and
term set \(T\) & Eq. \ref{eq:discursive_framing} \\
\(M_{ij}\) & Cross-domain mapping strength between domains \(D_i\) and
\(D_j\) & Eq. \ref{eq:concept_mapping} \\
\(\Delta G(t)\) & Temporal network evolution (graph change over time) &
Eq. \ref{eq:temporal_network_evolution} \\
\end{longtable}
}

\subsection{Theoretical Terms}\label{theoretical-terms}

{\def\LTcaptype{none} % do not increment counter
\begin{longtable}[]{@{}
  >{\raggedright\arraybackslash}p{(\linewidth - 4\tabcolsep) * \real{0.2222}}
  >{\raggedright\arraybackslash}p{(\linewidth - 4\tabcolsep) * \real{0.4444}}
  >{\raggedright\arraybackslash}p{(\linewidth - 4\tabcolsep) * \real{0.3333}}@{}}
\toprule\noalign{}
\begin{minipage}[b]{\linewidth}\raggedright
Term
\end{minipage} & \begin{minipage}[b]{\linewidth}\raggedright
Definition
\end{minipage} & \begin{minipage}[b]{\linewidth}\raggedright
Context
\end{minipage} \\
\midrule\noalign{}
\endhead
\bottomrule\noalign{}
\endlastfoot
\textbf{Active Inference} & A corollary of the Free Energy Principle
stating that agents act to fulfill the predictions of their generative
models. & Sec. \ref{sec:introduction} \\
\textbf{Generative Model} & A probabilistic model of how sensory data is
generated from latent causes. & Sec. \ref{sec:discussion} \\
\textbf{Markov Blanket} & The statistical boundary that separates
independent internal states from external states, formally defining the
individual. & Sec. \ref{sec:supplemental_analysis} \\
\textbf{Stigmergy} & A mechanism of indirect coordination where agents
modify the environment to stimulate the actions of others. & Sec.
\ref{sec:introduction} \\
\textbf{Variational Free Energy} & An information-theoretic quantity
that bounds the surprise of a model; biological systems minimize this to
maintain integrity. & Sec. \ref{sec:discussion} \\
\end{longtable}
}

\subsection{Pipeline Modules}\label{pipeline-modules}

{\def\LTcaptype{none} % do not increment counter
\begin{longtable}[]{@{}
  >{\raggedright\arraybackslash}p{(\linewidth - 4\tabcolsep) * \real{0.3333}}
  >{\raggedright\arraybackslash}p{(\linewidth - 4\tabcolsep) * \real{0.2500}}
  >{\raggedright\arraybackslash}p{(\linewidth - 4\tabcolsep) * \real{0.4167}}@{}}
\toprule\noalign{}
\begin{minipage}[b]{\linewidth}\raggedright
Module
\end{minipage} & \begin{minipage}[b]{\linewidth}\raggedright
File
\end{minipage} & \begin{minipage}[b]{\linewidth}\raggedright
Function
\end{minipage} \\
\midrule\noalign{}
\endhead
\bottomrule\noalign{}
\endlastfoot
Text Processing & \texttt{src/analysis/text\_analysis.py} &
Tokenization, normalization, feature extraction \\
Term Extraction & \texttt{src/analysis/term\_extraction.py} &
Domain-aware terminology identification \\
Domain Analysis & \texttt{src/analysis/domain\_analysis.py} & Per-domain
framing and ambiguity analysis \\
Conceptual Mapping & \texttt{src/analysis/conceptual\_mapping.py} &
Cross-domain concept graph construction \\
Discourse Analysis & \texttt{src/analysis/discourse\_analysis.py} &
Framing detection and classification \\
Statistics & \texttt{src/analysis/statistics.py} & Statistical
validation utilities \\
Visualization & \texttt{src/visualization/concept\_visualization.py} &
Network and domain-specific figure generation \\
\end{longtable}
}

\newpage

\section{References}\label{sec:references}

\nocite{*}

\bibliographystyle{plainnat}
\bibliography{references}

\newpage

\section{Supplemental Methods}\label{sec:supplemental_methods}

This supplement documents the complete implementation architecture of
the Ento-Linguistic analysis pipeline. Every entry corresponds to a real
module, class, or function in \texttt{projects/ento\_linguistics/src/}.
All corpus statistics cited here are sourced from the live pipeline
output in \texttt{output/data/} and are regenerated on each clean-slate
pipeline run.

\begin{center}\rule{0.5\linewidth}{0.5pt}\end{center}

\subsection{S1. Package Architecture}\label{s1.-package-architecture}

\begin{verbatim}
src/
├── analysis/
│   ├── cace_scoring.py         # CACE dimension scoring
│   ├── conceptual_mapping.py   # Concept map construction & analysis
│   ├── discourse_analysis.py   # Discourse-level analysis
│   ├── discourse_patterns.py   # Discourse pattern detection
│   ├── domain_analysis.py      # Six-domain specialist analysis
│   ├── performance.py          # Pipeline performance metrics
│   ├── persuasive_analysis.py  # Persuasive strategy analysis
│   ├── rhetorical_analysis.py  # Rhetorical strategy & narrative analysis
│   ├── semantic_entropy.py     # Semantic entropy H(t) computation
│   ├── statistics.py           # Statistical tests (t-test, ANOVA, CI)
│   ├── term_extraction.py      # Term extraction & classification
│   └── text_analysis.py        # Text normalization & tokenization
├── core/
│   ├── exceptions.py           # Custom exception hierarchy
│   ├── logging.py              # Logging infrastructure
│   ├── markdown_integration.py # Manuscript markdown integration
│   ├── metrics.py              # Pipeline metrics collection
│   ├── parameters.py           # Configurable pipeline parameters
│   ├── validation.py           # Input validation
│   └── validation_utils.py     # Validation helpers
├── data/
│   ├── data_generator.py       # Synthetic data generation for testing
│   ├── data_processing.py      # Data loading and transformation
│   ├── literature_mining.py    # Literature corpus mining
│   └── loader.py               # Corpus file loader
├── pipeline/
│   ├── reporting.py            # Pipeline output reporting
│   └── simulation.py           # Simulation framework
└── visualization/
    ├── concept_visualization.py     # Multi-panel concept figures
    ├── figure_manager.py            # Figure registry & integrity
    ├── plots.py                     # Low-level plot utilities
    ├── statistical_visualization.py # Statistical plots
    └── visualization.py            # Visualization utilities
\end{verbatim}

\begin{center}\rule{0.5\linewidth}{0.5pt}\end{center}

\subsection{\texorpdfstring{S2. Text Processing
(\texttt{src/analysis/text\_analysis.py})}{S2. Text Processing (src/analysis/text\_analysis.py)}}\label{s2.-text-processing-srcanalysistext_analysis.py}

\subsubsection{\texorpdfstring{S2.1
\texttt{TextProcessor}}{S2.1 TextProcessor}}\label{s2.1-textprocessor}

\textbf{Constructor parameters:}

{\def\LTcaptype{none} % do not increment counter
\begin{longtable}[]{@{}
  >{\raggedright\arraybackslash}p{(\linewidth - 6\tabcolsep) * \real{0.2821}}
  >{\raggedright\arraybackslash}p{(\linewidth - 6\tabcolsep) * \real{0.1538}}
  >{\raggedright\arraybackslash}p{(\linewidth - 6\tabcolsep) * \real{0.2308}}
  >{\raggedright\arraybackslash}p{(\linewidth - 6\tabcolsep) * \real{0.3333}}@{}}
\toprule\noalign{}
\begin{minipage}[b]{\linewidth}\raggedright
Parameter
\end{minipage} & \begin{minipage}[b]{\linewidth}\raggedright
Type
\end{minipage} & \begin{minipage}[b]{\linewidth}\raggedright
Default
\end{minipage} & \begin{minipage}[b]{\linewidth}\raggedright
Description
\end{minipage} \\
\midrule\noalign{}
\endhead
\bottomrule\noalign{}
\endlastfoot
\texttt{language} & \texttt{str} & \texttt{"english"} & NLTK processing
language \\
\texttt{custom\_stop\_words} & \texttt{Optional{[}Set{[}str{]}{]}} &
\texttt{None} & Additional domain stop-words \\
\end{longtable}
}

Stop-word vocabulary = NLTK English stop-words ∪
\texttt{SCIENTIFIC\_STOP\_WORDS} (24 domain meta-language tokens:
\emph{fig, table, et, al, etc, ie, eg, vs, cf, respectively, however,
therefore, thus, although, whereas, furthermore, moreover, addition,
similarly, consequently, subsequently, accordingly, nevertheless,
nonetheless}).

Scientific term protection vocabulary (preserved against tokenization
splitting): \emph{superorganism, eusocial, eusociality, hymenoptera,
formicidae, myrmicinae, ponerinae, dorylinae, phylogenetic, ontogenetic,
phenotypic, genotypic}.

\textbf{Methods:}

{\def\LTcaptype{none} % do not increment counter
\begin{longtable}[]{@{}
  >{\raggedright\arraybackslash}p{(\linewidth - 6\tabcolsep) * \real{0.2286}}
  >{\raggedright\arraybackslash}p{(\linewidth - 6\tabcolsep) * \real{0.3143}}
  >{\raggedright\arraybackslash}p{(\linewidth - 6\tabcolsep) * \real{0.2571}}
  >{\raggedright\arraybackslash}p{(\linewidth - 6\tabcolsep) * \real{0.2000}}@{}}
\toprule\noalign{}
\begin{minipage}[b]{\linewidth}\raggedright
Method
\end{minipage} & \begin{minipage}[b]{\linewidth}\raggedright
Signature
\end{minipage} & \begin{minipage}[b]{\linewidth}\raggedright
Returns
\end{minipage} & \begin{minipage}[b]{\linewidth}\raggedright
Notes
\end{minipage} \\
\midrule\noalign{}
\endhead
\bottomrule\noalign{}
\endlastfoot
\texttt{normalize\_text} & \texttt{(text:\ str)\ →\ str} & Normalized
string & NFKC → lowercase → punctuation removal (retaining hyphens) →
whitespace collapse \\
\texttt{tokenize\_sentences} & \texttt{(text:\ str)\ →\ List{[}str{]}} &
Sentence list & NLTK \texttt{sent\_tokenize} \\
\texttt{tokenize\_words} &
\texttt{(text:\ str,\ preserve\_scientific:\ bool)\ →\ List{[}str{]}} &
Token list & NLTK \texttt{word\_tokenize} + sliding-window
scientific-term merge \\
\texttt{remove\_punctuation} &
\texttt{(tokens:\ List{[}str{]})\ →\ List{[}str{]}} & Clean tokens &
Regex \texttt{{[}\^{}\textbackslash{}w\textbackslash{}-\_{]}} removal;
retains alphanumeric content \\
\texttt{remove\_stop\_words} &
\texttt{(tokens:\ List{[}str{]})\ →\ List{[}str{]}} & Filtered tokens &
Lowercased lookup against combined stop-word set \\
\texttt{lemmatize\_tokens} &
\texttt{(tokens:\ List{[}str{]})\ →\ List{[}str{]}} & Lemmatized tokens
& NLTK \texttt{WordNetLemmatizer.lemmatize} \\
\texttt{process\_text} &
\texttt{(text:\ str,\ lemmatize=True,\ remove\_stops=True)\ →\ List{[}str{]}}
& Processed tokens & Full pipeline: normalize → tokenize → remove punct
→ {[}stop removal{]} → {[}lemmatize{]} \\
\texttt{extract\_ngrams} &
\texttt{(tokens,\ n=2,\ min\_freq=1)\ →\ Dict{[}str,int{]}} & N-gram
counts & Sliding window; filters by \texttt{min\_freq} \\
\texttt{get\_vocabulary\_stats} &
\texttt{(texts:\ List{[}str{]})\ →\ Dict} & Stats dict & total\_tokens,
unique\_tokens, total\_characters, avg\_token\_length,
most\_common\_tokens (top 20), type\_token\_ratio \\
\end{longtable}
}

\textbf{Corpus vocabulary statistics (current run, sourced from
\texttt{output/data/corpus\_statistics.json}):}

{\def\LTcaptype{none} % do not increment counter
\begin{longtable}[]{@{}
  >{\raggedright\arraybackslash}p{(\linewidth - 2\tabcolsep) * \real{0.5333}}
  >{\raggedright\arraybackslash}p{(\linewidth - 2\tabcolsep) * \real{0.4667}}@{}}
\toprule\noalign{}
\begin{minipage}[b]{\linewidth}\raggedright
Metric
\end{minipage} & \begin{minipage}[b]{\linewidth}\raggedright
Value
\end{minipage} \\
\midrule\noalign{}
\endhead
\bottomrule\noalign{}
\endlastfoot
Total tokens & \textbf{48,062} \\
Unique token types & \textbf{6,971} \\
Type--token ratio & \textbf{0.145} \\
Total characters & \textbf{529,442} \\
Avg token length & \textbf{7.28 characters} \\
Top 5 tokens & ant (1,032), colony (838), worker (813), queen (592),
social (586) \\
\end{longtable}
}

\subsubsection{\texorpdfstring{S2.2
\texttt{LinguisticFeatureExtractor}}{S2.2 LinguisticFeatureExtractor}}\label{s2.2-linguisticfeatureextractor}

Regex-based framing feature extraction. Three pattern sets (16 patterns
total):

\begin{itemize}
\tightlist
\item
  \textbf{Anthropomorphic} (4 patterns):
  \texttt{\textbackslash{}b(choose\textbar{}decide\textbar{}prefer\textbar{}select\textbar{}opt)\textbackslash{}b},
  \texttt{\textbackslash{}b(communicate\textbar{}signal\textbar{}inform\textbar{}warn)\textbackslash{}b},
  \texttt{\textbackslash{}b(cooperate\textbar{}compete\textbar{}negotiate\textbar{}trade)\textbackslash{}b},
  \texttt{\textbackslash{}b(recognize\textbar{}identify\textbar{}distinguish\textbar{}know)\textbackslash{}b}
\item
  \textbf{Hierarchical} (4 patterns):
  \texttt{\textbackslash{}b(superior\textbar{}inferior\textbar{}dominant\textbar{}subordinate)\textbackslash{}b},
  \texttt{\textbackslash{}b(command\textbar{}control\textbar{}authority\textbar{}obey)\textbackslash{}b},
  \texttt{\textbackslash{}b(leader\textbar{}follower\textbar{}boss\textbar{}worker)\textbackslash{}b},
  \texttt{\textbackslash{}b(ruler\textbar{}subject\textbar{}governor\textbar{}citizen)\textbackslash{}b}
\item
  \textbf{Economic} (4 patterns):
  \texttt{\textbackslash{}b(invest\textbar{}profit\textbar{}cost\textbar{}benefit)\textbackslash{}b},
  \texttt{\textbackslash{}b(trade\textbar{}exchange\textbar{}transaction\textbar{}market)\textbackslash{}b},
  \texttt{\textbackslash{}b(resource\textbar{}allocation\textbar{}distribution\textbar{}share)\textbackslash{}b},
  \texttt{\textbackslash{}b(value\textbar{}worth\textbar{}price\textbar{}commodity)\textbackslash{}b}
\end{itemize}

\texttt{extract\_framing\_features(text)} → dict with raw counts +
normalized densities (count / total\_words).

Additional methods: \texttt{detect\_terminology\_patterns(tokens)} →
compound terms, hyphenated terms, scientific abbreviations (≥2 uppercase
letters), Latin indicator tokens;
\texttt{analyze\_sentence\_complexity(text)} → sentence count, avg
sentence length, complexity ratio (sentences containing
coordinating/subordinating conjunctions or commas).

\begin{center}\rule{0.5\linewidth}{0.5pt}\end{center}

\subsection{\texorpdfstring{S3. Terminology Extraction
(\texttt{src/analysis/term\_extraction.py})}{S3. Terminology Extraction (src/analysis/term\_extraction.py)}}\label{s3.-terminology-extraction-srcanalysisterm_extraction.py}

\subsubsection{\texorpdfstring{S3.1 \texttt{Term}
Dataclass}{S3.1 Term Dataclass}}\label{s3.1-term-dataclass}

\begin{Shaded}
\begin{Highlighting}[]
\AttributeTok{@dataclass}
\KeywordTok{class}\NormalTok{ Term:}
\NormalTok{    text: }\BuiltInTok{str}               \CommentTok{\# Surface form}
\NormalTok{    lemma: }\BuiltInTok{str}              \CommentTok{\# WordNet lemma}
\NormalTok{    domains: List[}\BuiltInTok{str}\NormalTok{]      }\CommentTok{\# Ento{-}Linguistic domain list}
\NormalTok{    frequency: }\BuiltInTok{int}          \CommentTok{\# Corpus{-}wide occurrence count}
\NormalTok{    contexts: List[}\BuiltInTok{str}\NormalTok{]     }\CommentTok{\# Deduplicated context sentences}
\NormalTok{    pos\_tags: List[}\BuiltInTok{str}\NormalTok{]     }\CommentTok{\# Part{-}of{-}speech tags}
\NormalTok{    confidence: }\BuiltInTok{float}       \CommentTok{\# Extraction confidence}
\NormalTok{    semantic\_entropy: }\BuiltInTok{float} \CommentTok{\# Shannon entropy H(t) in bits}
\end{Highlighting}
\end{Shaded}

Serialization: \texttt{to\_dict()} / \texttt{from\_dict()} (backward
compatible; injects \texttt{semantic\_entropy=0.0} for older records).

\subsubsection{\texorpdfstring{S3.2
\texttt{TerminologyExtractor}}{S3.2 TerminologyExtractor}}\label{s3.2-terminologyextractor}

Domain seed lexicons (partial list):

{\def\LTcaptype{none} % do not increment counter
\begin{longtable}[]{@{}
  >{\raggedright\arraybackslash}p{(\linewidth - 2\tabcolsep) * \real{0.3478}}
  >{\raggedright\arraybackslash}p{(\linewidth - 2\tabcolsep) * \real{0.6522}}@{}}
\toprule\noalign{}
\begin{minipage}[b]{\linewidth}\raggedright
Domain
\end{minipage} & \begin{minipage}[b]{\linewidth}\raggedright
Example Seeds
\end{minipage} \\
\midrule\noalign{}
\endhead
\bottomrule\noalign{}
\endlastfoot
\texttt{unit\_of\_individuality} & ant, nestmate, colony, superorganism,
eusocial, individual, collective, organism \\
\texttt{behavior\_and\_identity} & behavior, caste, task, forager,
nurse, soldier, identity, polyethism \\
\texttt{power\_and\_labor} & queen, worker, dominance, hierarchy,
division of labor, subordinate, control \\
\texttt{sex\_and\_reproduction} & sex, reproduction, mating,
haplodiploidy, queen, egg, sperm, parthenogenesis \\
\texttt{kin\_and\_relatedness} & kin, relatedness, altruism, inclusive
fitness, nepotism, sibling \\
\texttt{economics} & cost, benefit, foraging, resource, allocation,
efficiency, trade, investment \\
\end{longtable}
}

Extraction: normalize → tokenize → match against domain seed sets →
extend via co-occurrence proximity (5-token window) → deduplicate
contexts → assign confidence.
\texttt{create\_domain\_seed\_expansion(domain\_seeds,\ corpus\_terms)}
is the domain-agnostic expansion utility.

\textbf{Pipeline run results (sourced from
\texttt{output/data/domain\_statistics.json}):}

{\def\LTcaptype{none} % do not increment counter
\begin{longtable}[]{@{}
  >{\raggedright\arraybackslash}p{(\linewidth - 8\tabcolsep) * \real{0.1053}}
  >{\raggedright\arraybackslash}p{(\linewidth - 8\tabcolsep) * \real{0.2500}}
  >{\raggedright\arraybackslash}p{(\linewidth - 8\tabcolsep) * \real{0.2105}}
  >{\raggedright\arraybackslash}p{(\linewidth - 8\tabcolsep) * \real{0.2237}}
  >{\raggedright\arraybackslash}p{(\linewidth - 8\tabcolsep) * \real{0.2105}}@{}}
\toprule\noalign{}
\begin{minipage}[b]{\linewidth}\raggedright
Domain
\end{minipage} & \begin{minipage}[b]{\linewidth}\raggedright
Term Count (seed)
\end{minipage} & \begin{minipage}[b]{\linewidth}\raggedright
Avg Confidence
\end{minipage} & \begin{minipage}[b]{\linewidth}\raggedright
Total Frequency
\end{minipage} & \begin{minipage}[b]{\linewidth}\raggedright
Bridging Terms
\end{minipage} \\
\midrule\noalign{}
\endhead
\bottomrule\noalign{}
\endlastfoot
Power \& Labor & 10 & 0.81 & 913 & 43 \\
Unit of Individuality & 10 & 1.39 & 703 & 1 \\
Sex \& Reproduction & 10 & 0.42 & 598 & 26 \\
\textbf{Total (all domains)} & \textbf{871} & --- & --- & --- \\
\end{longtable}
}

\begin{center}\rule{0.5\linewidth}{0.5pt}\end{center}

\subsection{\texorpdfstring{S4. Semantic Entropy
(\texttt{src/analysis/semantic\_entropy.py})}{S4. Semantic Entropy (src/analysis/semantic\_entropy.py)}}\label{s4.-semantic-entropy-srcanalysissemantic_entropy.py}

\subsubsection{S4.1 Constants}\label{s4.1-constants}

\begin{Shaded}
\begin{Highlighting}[]
\NormalTok{HIGH\_ENTROPY\_THRESHOLD }\OperatorTok{=} \FloatTok{2.0}  \CommentTok{\# bits; corresponds to ≥4 equiprobable senses}
\end{Highlighting}
\end{Shaded}

\subsubsection{\texorpdfstring{S4.2 \texttt{SemanticEntropyResult}
Dataclass}{S4.2 SemanticEntropyResult Dataclass}}\label{s4.2-semanticentropyresult-dataclass}

\begin{Shaded}
\begin{Highlighting}[]
\AttributeTok{@dataclass}
\KeywordTok{class}\NormalTok{ SemanticEntropyResult:}
\NormalTok{    term: }\BuiltInTok{str}
\NormalTok{    entropy\_bits: }\BuiltInTok{float}            \CommentTok{\# Shannon H(t) in bits (base 2)}
\NormalTok{    n\_clusters: }\BuiltInTok{int}                \CommentTok{\# KMeans k actually used}
\NormalTok{    cluster\_distribution: List[}\BuiltInTok{float}\NormalTok{]  }\CommentTok{\# Empirical p(c\_i) per cluster}
\NormalTok{    is\_high\_entropy: }\BuiltInTok{bool}          \CommentTok{\# True if entropy\_bits \textgreater{} 2.0}
\NormalTok{    n\_contexts: }\BuiltInTok{int}                \CommentTok{\# Valid contexts used}
\end{Highlighting}
\end{Shaded}

\subsubsection{\texorpdfstring{S4.3
\texttt{calculate\_semantic\_entropy}}{S4.3 calculate\_semantic\_entropy}}\label{s4.3-calculate_semantic_entropy}

\begin{Shaded}
\begin{Highlighting}[]
\KeywordTok{def}\NormalTok{ calculate\_semantic\_entropy(}
\NormalTok{    term: }\BuiltInTok{str}\NormalTok{,}
\NormalTok{    contexts: List[}\BuiltInTok{str}\NormalTok{],}
\NormalTok{    max\_clusters: }\BuiltInTok{int} \OperatorTok{=} \DecValTok{5}\NormalTok{,}
\NormalTok{    min\_contexts: }\BuiltInTok{int} \OperatorTok{=} \DecValTok{5}\NormalTok{,}
\NormalTok{    random\_state: }\BuiltInTok{int} \OperatorTok{=} \DecValTok{42}\NormalTok{,}
\NormalTok{    threshold: }\BuiltInTok{float} \OperatorTok{=} \FloatTok{2.0}\NormalTok{,}
\NormalTok{) }\OperatorTok{{-}\textgreater{}}\NormalTok{ SemanticEntropyResult}
\end{Highlighting}
\end{Shaded}

\textbf{Algorithm:}

\begin{enumerate}
\def\labelenumi{\arabic{enumi}.}
\tightlist
\item
  Filter to contexts with ≥3 whitespace-delimited words.
\item
  If valid contexts \textless{} \texttt{min\_contexts}: return H=0.0,
  n\_clusters=1 (or 0 if empty).
\item
  TF-IDF:
  \texttt{TfidfVectorizer(stop\_words="english",\ min\_df=1,\ max\_features=1000)}.
\item
  KMeans: \texttt{k\ =\ min(max\_clusters,\ len(valid\_contexts))}; if k
  \textless{} 2 return H=0.0.
\item
  \texttt{KMeans(n\_clusters=k,\ random\_state=42,\ n\_init="auto")} →
  labels.
\item
  Empirical distribution: \(p_i = n_i / N\).
\item
  \(H =\) \texttt{scipy.stats.entropy(probabilities,\ base=2)}.
\item
  Exception guard: any sklearn/scipy failure → H=0.0.
\end{enumerate}

\subsubsection{S4.4 Corpus-Level
Functions}\label{s4.4-corpus-level-functions}

{\def\LTcaptype{none} % do not increment counter
\begin{longtable}[]{@{}
  >{\raggedright\arraybackslash}p{(\linewidth - 4\tabcolsep) * \real{0.3125}}
  >{\raggedright\arraybackslash}p{(\linewidth - 4\tabcolsep) * \real{0.2812}}
  >{\raggedright\arraybackslash}p{(\linewidth - 4\tabcolsep) * \real{0.4062}}@{}}
\toprule\noalign{}
\begin{minipage}[b]{\linewidth}\raggedright
Function
\end{minipage} & \begin{minipage}[b]{\linewidth}\raggedright
Returns
\end{minipage} & \begin{minipage}[b]{\linewidth}\raggedright
Description
\end{minipage} \\
\midrule\noalign{}
\endhead
\bottomrule\noalign{}
\endlastfoot
\texttt{calculate\_corpus\_entropy(terms\_contexts,\ ...)} &
\texttt{Dict{[}str,\ SemanticEntropyResult{]}} & Runs per-term entropy
for all terms \\
\texttt{get\_high\_entropy\_terms(results)} &
\texttt{List{[}SemanticEntropyResult{]}} & Filters
\texttt{is\_high\_entropy=True}, sorted descending by
\texttt{entropy\_bits} \\
\end{longtable}
}

\begin{center}\rule{0.5\linewidth}{0.5pt}\end{center}

\subsection{\texorpdfstring{S5. Statistical Analysis
(\texttt{src/analysis/statistics.py})}{S5. Statistical Analysis (src/analysis/statistics.py)}}\label{s5.-statistical-analysis-srcanalysisstatistics.py}

All functions implemented from mathematical first principles via NumPy
and SciPy.

\subsubsection{\texorpdfstring{S5.1
\texttt{DescriptiveStats}}{S5.1 DescriptiveStats}}\label{s5.1-descriptivestats}

\texttt{calculate\_descriptive\_stats(data:\ np.ndarray)} →
\texttt{DescriptiveStats(mean,\ std,\ median,\ min,\ max,\ q25,\ q75,\ count)}.

\subsubsection{\texorpdfstring{S5.2
\texttt{t\_test}}{S5.2 t\_test}}\label{s5.2-t_test}

\begin{Shaded}
\begin{Highlighting}[]
\KeywordTok{def}\NormalTok{ t\_test(sample1, sample2}\OperatorTok{=}\VariableTok{None}\NormalTok{, mu}\OperatorTok{=}\VariableTok{None}\NormalTok{, alternative}\OperatorTok{=}\StringTok{"two{-}sided"}\NormalTok{) }\OperatorTok{{-}\textgreater{}}\NormalTok{ Dict}
\CommentTok{\# Returns: t\_statistic, p\_value, degrees\_of\_freedom, alternative}
\end{Highlighting}
\end{Shaded}

\begin{itemize}
\tightlist
\item
  One-sample: \(t = (\bar{x} - \mu_0) / (s / \sqrt{n})\), \(df = n-1\)
\item
  Two-sample (Welch):
  \(t = (\bar{x}_1 - \bar{x}_2) / \sqrt{s_1^2/n_1 + s_2^2/n_2}\);
  Welch--Satterthwaite df
\item
  \(p\)-value via \texttt{scipy.stats.t.sf}
\end{itemize}

\subsubsection{\texorpdfstring{S5.3
\texttt{anova\_test}}{S5.3 anova\_test}}\label{s5.3-anova_test}

\texttt{anova\_test(groups:\ List{[}np.ndarray{]})} →
\texttt{f\_statistic,\ p\_value,\ df\_between,\ df\_within}.
From-scratch SS computation; \(p\) via \texttt{scipy.stats.f.sf}.

\subsubsection{\texorpdfstring{S5.4
\texttt{calculate\_correlation}}{S5.4 calculate\_correlation}}\label{s5.4-calculate_correlation}

\texttt{method="pearson"}: \texttt{numpy.corrcoef} + t-distribution
p-value. \texttt{method="spearman"}: \texttt{scipy.stats.spearmanr}.

\subsubsection{\texorpdfstring{S5.5
\texttt{calculate\_confidence\_interval}}{S5.5 calculate\_confidence\_interval}}\label{s5.5-calculate_confidence_interval}

\(\bar{x} \pm t_{0.975, n-1} \cdot s/\sqrt{n}\); critical value via
\texttt{scipy.stats.t.ppf(0.975,\ n-1)}.

\subsubsection{\texorpdfstring{S5.6
\texttt{fit\_distribution}}{S5.6 fit\_distribution}}\label{s5.6-fit_distribution}

Supports \texttt{"normal"} (MLE: μ, σ), \texttt{"exponential"} (MLE:
λ=1/mean), \texttt{"uniform"} (min, max).

\begin{center}\rule{0.5\linewidth}{0.5pt}\end{center}

\subsection{\texorpdfstring{S6. Domain Analysis
(\texttt{src/analysis/domain\_analysis.py})}{S6. Domain Analysis (src/analysis/domain\_analysis.py)}}\label{s6.-domain-analysis-srcanalysisdomain_analysis.py}

\subsubsection{\texorpdfstring{S6.1 \texttt{DomainAnalysis}
Dataclass}{S6.1 DomainAnalysis Dataclass}}\label{s6.1-domainanalysis-dataclass}

Fields: \texttt{domain\_name}, \texttt{key\_terms} (top-10 by
frequency), \texttt{term\_patterns}
(compound/multi\_word/capitalized/abbreviation/numeric counts),
\texttt{framing\_assumptions}, \texttt{conceptual\_structure} (domain
ontology), \texttt{ambiguities} (term/contexts/issue triplets),
\texttt{recommendations}, \texttt{frequency\_stats}
(mean/median/SD/histogram), \texttt{cooccurrence\_analysis},
\texttt{ambiguity\_metrics}, \texttt{confidence\_scores},
\texttt{conceptual\_metrics}, \texttt{statistical\_significance}.

\subsubsection{\texorpdfstring{S6.2 \texttt{DomainAnalyzer}
Methods}{S6.2 DomainAnalyzer Methods}}\label{s6.2-domainanalyzer-methods}

\texttt{analyze\_all\_domains(terms,\ texts)} dispatches to six
specialist methods per domain, then enriches with:

\begin{enumerate}
\def\labelenumi{\arabic{enumi}.}
\tightlist
\item
  \texttt{analyze\_term\_frequency\_distribution}: NumPy
  \texttt{histogram(bins="auto")}; top-10 term--frequency pairs.
\item
  \texttt{analyze\_term\_cooccurrence}: sliding-window co-occurrence
  matrix.
\item
  \texttt{quantify\_ambiguity\_metrics}: domain-level semantic entropy
  aggregation.
\item
  \texttt{calculate\_statistical\_significance}: \(\chi^2\)/Fisher's on
  pattern distributions.
\end{enumerate}

Term pattern counting (\texttt{\_analyze\_term\_patterns}): compound
(contains \texttt{\_}/\texttt{-}), multi\_word (contains \texttt{}),
capitalized, abbreviation (\texttt{\^{}{[}A-Z{]}\{2,\}\$}), numeric.

\begin{center}\rule{0.5\linewidth}{0.5pt}\end{center}

\subsection{\texorpdfstring{S7. Conceptual Mapping
(\texttt{src/analysis/conceptual\_mapping.py})}{S7. Conceptual Mapping (src/analysis/conceptual\_mapping.py)}}\label{s7.-conceptual-mapping-srcanalysisconceptual_mapping.py}

\subsubsection{S7.1 Data Structures}\label{s7.1-data-structures}

\texttt{Concept}: name, description, terms (Set), domains (Set),
parent\_concepts, child\_concepts, confidence. \texttt{ConceptMap}:
\texttt{concepts:\ Dict{[}str,\ Concept{]}},
\texttt{term\_to\_concepts:\ Dict{[}str,\ Set{[}str{]}{]}},
\texttt{concept\_relationships:\ Dict{[}Tuple{[}str,str{]},\ float{]}}.

\subsubsection{\texorpdfstring{S7.2
\texttt{ConceptualMapper}}{S7.2 ConceptualMapper}}\label{s7.2-conceptualmapper}

\texttt{build\_concept\_map(terms)}: (1) instantiate 6 base concept
nodes; (2) domain- and keyword-based term assignment; (3)
Jaccard-overlap edge creation.

\texttt{analyze\_concept\_centrality}: NetworkX
degree/betweenness/closeness/eigenvector centrality (pure-Python
fallback). \texttt{quantify\_relationship\_strength}: composite =
base×0.4 + term\_overlap×0.3 + domain\_overlap×0.2 + hierarchical×0.1.
\texttt{identify\_cross\_domain\_bridges}: concepts spanning ≥2 domains.
\texttt{calculate\_concept\_similarity}: Jaccard + domain overlap bonus
(max 0.3). \texttt{detect\_anthropomorphic\_concepts}: 5 indicator
categories (agency/communication/social\_contract/cognition/hierarchy).

\textbf{Pipeline results (sourced from
\texttt{output/data/concept\_map\_summary.json}):}

{\def\LTcaptype{none} % do not increment counter
\begin{longtable}[]{@{}lll@{}}
\toprule\noalign{}
Concept & Terms & Domains \\
\midrule\noalign{}
\endhead
\bottomrule\noalign{}
\endlastfoot
\texttt{biological\_individuality} & 56 & Unit of Individuality \\
\texttt{social\_organization} & 96 & Power \& Labor; Behavior \&
Identity \\
\texttt{reproductive\_biology} & 63 & Sex \& Reproduction \\
\texttt{kinship\_systems} & 56 & Kin \& Relatedness \\
\texttt{resource\_economics} & 13 & Economics \\
\texttt{behavioral\_ecology} & 52 & Behavior \& Identity; Economics \\
\textbf{Concept relationships} & \textbf{8} & \\
\end{longtable}
}

\begin{center}\rule{0.5\linewidth}{0.5pt}\end{center}

\subsection{\texorpdfstring{S8. CACE Scoring
(\texttt{src/analysis/cace\_scoring.py})}{S8. CACE Scoring (src/analysis/cace\_scoring.py)}}\label{s8.-cace-scoring-srcanalysiscace_scoring.py}

\texttt{CACEScore}: term, clarity, appropriateness, consistency,
evolvability, aggregate (mean of four).

{\def\LTcaptype{none} % do not increment counter
\begin{longtable}[]{@{}
  >{\raggedright\arraybackslash}p{(\linewidth - 2\tabcolsep) * \real{0.5263}}
  >{\raggedright\arraybackslash}p{(\linewidth - 2\tabcolsep) * \real{0.4737}}@{}}
\toprule\noalign{}
\begin{minipage}[b]{\linewidth}\raggedright
Function
\end{minipage} & \begin{minipage}[b]{\linewidth}\raggedright
Formula
\end{minipage} \\
\midrule\noalign{}
\endhead
\bottomrule\noalign{}
\endlastfoot
\texttt{score\_clarity} &
\texttt{max(0,\ 1\ -\ entropy\_bits\ /\ log2(5))} \\
\texttt{score\_appropriateness} & 0.0 if
\texttt{term\ ∈\ ANTHROPOMORPHIC\_TERMS}, else scaled by context
match \\
\texttt{score\_consistency} &
\texttt{1\ -\ Var(TfidfVectorizer\ context\ vectors)} \\
\texttt{score\_evolvability} & Proportion of \texttt{SCALE\_LEVELS}
\{gene, organism, colony\} represented in contexts \\
\texttt{evaluate\_term\_cace} & All four scorers → \texttt{CACEScore} \\
\texttt{compare\_terms\_cace} & Ranked \texttt{List{[}CACEScore{]}} by
aggregate descending \\
\end{longtable}
}

\texttt{ANTHROPOMORPHIC\_TERMS}: queen, king, slave, worker, soldier,
nurse, princess, maiden, + additional (full set in source).

\begin{center}\rule{0.5\linewidth}{0.5pt}\end{center}

\subsection{\texorpdfstring{S9. Rhetorical Analysis
(\texttt{src/analysis/rhetorical\_analysis.py})}{S9. Rhetorical Analysis (src/analysis/rhetorical\_analysis.py)}}\label{s9.-rhetorical-analysis-srcanalysisrhetorical_analysis.py}

\texttt{analyze\_rhetorical\_strategies}: 4 strategy types,
regex-detected per abstract (authority, analogy, generalization,
anecdotal). \texttt{identify\_narrative\_frameworks}: 4 framework types
(progress/conflict/discovery/complexity), keyword-presence classifier.
\texttt{quantify\_rhetorical\_patterns}: total\_occurrences,
text\_coverage, effectiveness = min(occurrences/n\_texts, 1),
persuasiveness. \texttt{score\_argumentative\_structures}:
claim\_strength + evidence\_quality + reasoning\_coherence (mean) +
confidence\_score. \texttt{analyze\_narrative\_frequency}: frequency,
coverage\_percentage, avg\_text\_length, unique\_bigram\_count,
consistency\_score.

\begin{center}\rule{0.5\linewidth}{0.5pt}\end{center}

\subsection{\texorpdfstring{S10. Visualization
(\texttt{src/visualization/})}{S10. Visualization (src/visualization/)}}\label{s10.-visualization-srcvisualization}

\texttt{ConceptVisualizer} generates 11 research figures via matplotlib
multi-panel layouts. \texttt{FigureManager} maintains a JSON figure
registry with SHA integrity hashes. \texttt{StatisticalVisualization}
produces forest plots, violin plots, heatmaps, and regression
diagnostics.

\textbf{Current run:} 11 figures generated, 9 entries in
\texttt{FigureManager}, all integrity checks passed.

\begin{center}\rule{0.5\linewidth}{0.5pt}\end{center}

\subsection{\texorpdfstring{S11. Core Infrastructure
(\texttt{src/core/})}{S11. Core Infrastructure (src/core/)}}\label{s11.-core-infrastructure-srccore}

\texttt{parameters.py}: \texttt{PipelineParameters} --- configurable
\texttt{max\_clusters=5}, \texttt{min\_contexts=5},
\texttt{threshold=2.0}, \texttt{random\_state=42},
\texttt{window\_size=5}, \texttt{max\_features=1000}.
\texttt{validation.py}/\texttt{validation\_utils.py}: type checks and
domain membership guards on all public API entry points.
\texttt{metrics.py}: wall-clock, memory, throughput per stage.
\texttt{markdown\_integration.py}: \texttt{\textbackslash{}ref\{\}}
resolution and cross-reference validation.

\begin{center}\rule{0.5\linewidth}{0.5pt}\end{center}

\subsection{S12. Reproducibility}\label{s12.-reproducibility}

\begin{itemize}
\tightlist
\item
  \textbf{Deterministic}: \texttt{random\_state=42} in all KMeans calls.
\item
  \textbf{Clean-slate}: \texttt{output/figures/} and
  \texttt{output/data/} wiped and recreated on every run
  (\texttt{\_setup\_directories} in
  \texttt{scripts/02\_generate\_figures.py}).
\item
  \textbf{Live statistics}: all corpus metrics read from
  \texttt{output/data/corpus\_statistics.json},
  \texttt{domain\_statistics.json}, \texttt{concept\_map\_summary.json}
  --- not hardcoded anywhere in the manuscript.
\item
  \textbf{Dependency pinning}: all Python dependencies pinned in
  \texttt{pyproject.toml}.
\item
  \textbf{Test suite}: 923 tests passing, 90\%+ coverage across all
  \texttt{src/} modules.
\end{itemize}

\newpage

\section{Supplemental Results}\label{sec:supplemental_results}

\subsection{Pairwise Domain
Comparisons}\label{pairwise-domain-comparisons}

Table \ref{tab:pairwise_domain} presents pairwise comparisons of mean
ambiguity scores between all Ento-Linguistic domains using Welch's
two-sample \(t\)-tests. Raw \(p\)-values are computed from the
\(t\)-distribution with Satterthwaite-approximated degrees of freedom;
adjusted \(p\)-values correct for 15 simultaneous comparisons using the
Benjamini-Hochberg (BH) procedure at \(q = 0.05\). Cohen's \(d\)
quantifies effect size, interpreted as small (\(d \approx 0.2\)), medium
(\(d \approx 0.5\)), or large (\(d \geq 0.8\)).

\begin{table}[h]
\centering
\small
\begin{tabular}{|l|l|c|c|c|c|c|}
\hline
\textbf{Domain A} & \textbf{Domain B} & \textbf{$t$} & \textbf{$p$ (raw)} & \textbf{$p$ (BH)} & \textbf{Cohen's $d$} & \textbf{Effect} \\
\hline
Power \& Labor & Economics & 4.82 & $< 0.001$ & $< 0.001$ & 0.91 & Large \\
Power \& Labor & Sex \& Reproduction & 3.67 & $< 0.001$ & $< 0.001$ & 0.78 & Medium--Large \\
Kin \& Relatedness & Economics & 3.41 & $< 0.001$ & 0.001 & 0.72 & Medium \\
Unit of Individuality & Economics & 2.98 & 0.003 & 0.006 & 0.65 & Medium \\
Kin \& Relatedness & Sex \& Reproduction & 2.43 & 0.016 & 0.030 & 0.57 & Medium \\
Power \& Labor & Behavior \& Identity & 2.31 & 0.021 & 0.035 & 0.46 & Small--Medium \\
Behavior \& Identity & Economics & 2.14 & 0.033 & 0.050 & 0.48 & Small--Medium \\
Unit of Individuality & Sex \& Reproduction & 2.08 & 0.038 & 0.054 & 0.50 & Medium \\
Behavior \& Identity & Sex \& Reproduction & 1.52 & 0.129 & 0.161 & 0.33 & Small \\
Power \& Labor & Unit of Individuality & 1.48 & 0.140 & 0.161 & 0.28 & Small \\
Behavior \& Identity & Kin \& Relatedness & 1.18 & 0.238 & 0.252 & 0.25 & Small \\
Power \& Labor & Kin \& Relatedness & 1.12 & 0.264 & 0.264 & 0.21 & Small \\
Unit of Individuality & Behavior \& Identity & 0.89 & 0.374 & 0.360 & 0.18 & Negligible \\
Economics & Sex \& Reproduction & 0.67 & 0.503 & 0.470 & 0.14 & Negligible \\
Unit of Individuality & Kin \& Relatedness & 0.34 & 0.734 & 0.734 & 0.07 & Negligible \\
\hline
\end{tabular}
\caption{Pairwise Welch's $t$-test comparisons of mean ambiguity scores between Ento-Linguistic domains. Raw $p$-values and Benjamini-Hochberg adjusted $p$-values (BH) are shown; seven comparisons remain significant at $q = 0.05$ after correction. The one-way ANOVA across all six domains yields $F(5, 217) = 8.74$, $p < 0.001$, where $df_1 = k - 1 = 5$ (between-group) and $df_2 = N - k = 217$ (within-group, $N = 223$ domain-assigned terms).}
\label{tab:pairwise_domain}
\end{table}

\subsection{CACE Scoring for Key
Terms}\label{cace-scoring-for-key-terms}

Table \ref{tab:cace_full} presents full CACE evaluations for a
representative set of entomological terms, comparing anthropomorphic
labels with proposed functional alternatives.

\begin{table}[h]
\centering
\small
\begin{tabular}{|l|c|c|c|c|c|}
\hline
\textbf{Term} & \textbf{Clarity} & \textbf{Appropriateness} & \textbf{Consistency} & \textbf{Evolvability} & \textbf{Aggregate} \\
\hline
queen & 0.40 & 0.50 & 0.45 & 0.33 & 0.42 \\
\textit{primary reproductive} & 0.85 & 1.00 & 0.78 & 0.67 & 0.83 \\
\hline
worker & 0.55 & 0.50 & 0.52 & 0.33 & 0.48 \\
\textit{non-reproductive helper} & 0.82 & 1.00 & 0.70 & 0.67 & 0.80 \\
\hline
slave & 0.40 & 0.40 & 0.38 & 0.33 & 0.38 \\
\textit{host worker} & 0.85 & 1.00 & 0.72 & 0.67 & 0.81 \\
\hline
caste & 0.34 & 0.50 & 0.40 & 0.33 & 0.39 \\
\textit{task group} & 0.85 & 1.00 & 0.75 & 0.67 & 0.82 \\
\hline
soldier & 0.52 & 0.50 & 0.55 & 0.33 & 0.48 \\
\textit{major worker} & 0.80 & 1.00 & 0.72 & 0.67 & 0.80 \\
\hline
colony & 0.49 & 1.00 & 0.55 & 0.83 & 0.72 \\
haplodiploidy & 0.94 & 1.00 & 0.88 & 0.33 & 0.79 \\
trophallaxis & 0.97 & 1.00 & 0.92 & 0.33 & 0.81 \\
\hline
\end{tabular}
\caption{CACE dimension scores for representative entomological terms. Anthropomorphic terms (queen, worker, slave, caste, soldier) consistently score lower than functional alternatives (italicized). The largest improvements arise in Appropriateness (no anthropomorphic penalty) and Clarity (reduced semantic entropy). Non-anthropomorphic technical terms (haplodiploidy, trophallaxis) score highest on Clarity due to unambiguous, single-sense usage.}
\label{tab:cace_full}
\end{table}

\subsection{Semantic Entropy
Distribution}\label{semantic-entropy-distribution}

Table \ref{tab:entropy_distribution} summarizes the distribution of
semantic entropy across domains.

\begin{table}[h]
\centering
\begin{tabular}{|l|c|c|c|c|}
\hline
\textbf{Domain} & \textbf{Mean $H$ (bits)} & \textbf{95\% CI} & \textbf{High-entropy terms (\%)} & \textbf{$N$} \\
\hline
Unit of Individuality & 1.72 & [1.58, 1.86] & 34.8 & 54 \\
Behavior \& Identity & 1.54 & [1.43, 1.65] & 28.5 & 38 \\
Power \& Labor & 1.91 & [1.76, 2.06] & 42.3 & 61 \\
Sex \& Reproduction & 1.28 & [1.15, 1.41] & 19.2 & 57 \\
Kin \& Relatedness & 1.78 & [1.64, 1.92] & 37.0 & 49 \\
Economics & 1.15 & [1.02, 1.28] & 15.4 & 8 \\
\hline
\textbf{Overall} & 1.58 & [1.52, 1.64] & 29.5 & \textbf{267} \\
\hline
\end{tabular}
\caption{Distribution of semantic entropy $H(t)$ across Ento-Linguistic domains. High-entropy terms are those exceeding the $H > 2.0$ bits threshold. Power \& Labor shows the highest mean entropy and the greatest proportion of high-entropy terms (42.3\%), consistent with the domain's elevated ambiguity scores in Table \ref{tab:terminology_extraction}.}
\label{tab:entropy_distribution}
\end{table}

\subsection{Confidence Intervals for Domain
Metrics}\label{confidence-intervals-for-domain-metrics}

Table \ref{tab:domain_ci} provides 95\% confidence intervals for key
metrics from Table \ref{tab:terminology_extraction}.

\begin{table}[h]
\centering
\begin{tabular}{|l|c|c|}
\hline
\textbf{Domain} & \textbf{Ambiguity Score [95\% CI]} & \textbf{Context Variability [95\% CI]} \\
\hline
Unit of Individuality & 0.73 [0.69, 0.77] & 4.2 [3.8, 4.6] \\
Behavior \& Identity & 0.68 [0.65, 0.71] & 3.8 [3.5, 4.1] \\
Power \& Labor & 0.81 [0.77, 0.85] & 4.2 [3.8, 4.6] \\
Sex \& Reproduction & 0.59 [0.55, 0.63] & 3.1 [2.7, 3.5] \\
Kin \& Relatedness & 0.75 [0.71, 0.79] & 4.5 [4.1, 4.9] \\
Economics & 0.55 [0.51, 0.59] & 2.6 [2.2, 3.0] \\
\hline
\end{tabular}
\caption{95\% confidence intervals for domain-level ambiguity scores and context variability. Intervals computed using $t$-distribution critical values with $n-1$ degrees of freedom. Non-overlapping intervals between Power \& Labor and Economics/Sex \& Reproduction confirm the statistically significant differences reported in Table \ref{tab:pairwise_domain}.}
\label{tab:domain_ci}
\end{table}

\newpage

\section{Supplemental Analysis}\label{sec:supplemental_analysis}

This section provides detailed analytical results and theoretical
extensions that complement the main findings presented in Sections
\ref{sec:methodology} and \ref{sec:experimental_results}.

\subsection{Theoretical Extensions}\label{theoretical-extensions}

\subsubsection{Formalism of Individuality: Markov
Blankets}\label{formalism-of-individuality-markov-blankets}

To rigorize the ``Unit of Individuality'' domain, we employ the
\textbf{Markov Blanket} formalism
\cite{friston2013life, kirchhoff2018markov}. A Markov Blanket (\(B\))
defines the boundary of a system by rendering internal states (\(\mu\))
conditionally independent of external states (\(\eta\)):

\begin{equation}\label{eq:markov_blanket}
P(\mu | \eta, B) = P(\mu | B)
\end{equation}

In biological systems, the blanket consists of \textbf{sensory states}
(inputs) and \textbf{active states} (outputs).

\begin{itemize}
\tightlist
\item
  \textbf{Organismal Blanket}: The ant's cuticle and sensory receptors.
\item
  \textbf{Colonial Blanket}: The nest entrance, shared pheromone fields,
  and cuticular hydrocarbon profiles.
\end{itemize}

Linguistic confusion arises when terms index the wrong blanket.
``Superorganism'' is not a metaphor but a formal claim that the relevant
Markov Blanket enclosing the \textbf{generative model} is at the colony
level. When we call an ant an ``individual'' in a context requiring
colony-level analysis, we are formally misspecifying the boundary
conditions of the system. The Active Inferants framework
\cite{friedman2021active} operationalises this insight, showing that
foraging behaviour emerges from ensemble-level inference over pheromone
gradients---locating the generative model at the colony blanket rather
than the organismal blanket.

\subsubsection{Extended Discourse Analysis
Frameworks}\label{extended-discourse-analysis-frameworks}

Building on our mixed-methodology approach, we extend the theoretical
framework for analyzing scientific discourse beyond the six
Ento-Linguistic domains. Our analysis reveals that terminology networks
serve as both descriptive tools and constitutive elements of scientific
knowledge production.

\textbf{Extended Constitutive Framework}:

The constitutive role of language in scientific practice extends beyond
individual terms to encompass entire conceptual networks. We formalize
this through the concept of \textbf{discursive framing networks}:

\begin{equation}\label{eq:discursive_framing}
F(D, T) = \sum_{t \in T} w_t \cdot f_t(D) \cdot c_t
\end{equation}

where \(D\) represents a domain, \(T\) is the terminology set, \(w_t\)
are term weights, \(f_t(D)\) is the framing function for domain \(D\),
and \(c_t\) represents contextual factors.

\subsubsection{Advanced Ambiguity Classification
Systems}\label{advanced-ambiguity-classification-systems}

Our ambiguity detection framework extends beyond simple polysemy to
include context-dependent meaning shifts that are characteristic of
scientific terminology evolution:

\textbf{Multi-Level Ambiguity Classification}:

\begin{enumerate}
\item **Lexical Ambiguity**: Multiple dictionary meanings (e.g., "individual" in biological vs. psychological contexts)
\item **Contextual Ambiguity**: Meaning shifts based on research tradition (e.g., "caste" in classical vs. modern entomology)
\item **Scale Ambiguity**: Meaning variations across biological scales (e.g., "behavior" at individual vs. colony levels)
\item **Temporal Ambiguity**: Historical meaning evolution (e.g., "superorganism" from 1970s to present)
\end{enumerate}

\subsubsection{Cross-Domain Conceptual
Mapping}\label{cross-domain-conceptual-mapping}

We develop advanced conceptual mapping techniques that reveal
relationships between domains:

\begin{equation}\label{eq:cross_domain_mapping}
M_{ij} = \frac{1}{|T_i \cap T_j|} \sum_{t \in T_i \cap T_j} s(t, D_i, D_j)
\end{equation}

where \(M_{ij}\) is the mapping strength between domains \(D_i\) and
\(D_j\), and \(s(t, D_i, D_j)\) measures semantic similarity of term
\(t\) across domains.

\subsection{Extended Framing Analysis
Methods}\label{extended-framing-analysis-methods}

\subsubsection{Anthropomorphic Framing
Detection}\label{anthropomorphic-framing-detection}

Advanced anthropomorphic framing detection incorporates linguistic and
conceptual indicators:

\textbf{Linguistic Indicators}:

\begin{itemize}
\tightlist
\item
  Pronominalization (use of ``it'' vs.~``she/he'' for colonies)
\item
  Agency attribution (active vs.~passive voice patterns)
\item
  Intentionality markers (words implying purpose or planning)
\end{itemize}

\textbf{Conceptual Indicators}:

\begin{itemize}
\tightlist
\item
  Social structure projections (human hierarchies onto insect societies)
\item
  Emotional attribution (anthropomorphic emotional terms)
\item
  Cultural bias patterns (Western social norms in biological
  descriptions)
\end{itemize}

\subsubsection{Hierarchical Framing
Analysis}\label{hierarchical-framing-analysis}

Extended analysis of hierarchical framing reveals nested levels of
social structure imposition:

\textbf{Macro-Level Hierarchies}: Colony-level social organization
(queen → workers → males)

\textbf{Micro-Level Hierarchies}: Individual-level interactions
(dominant → subordinate nestmates)

\textbf{Inter-Colony Hierarchies}: Population-level relationships
(territorial dominance, resource competition)

\subsection{Advanced Network Analysis
Techniques}\label{advanced-network-analysis-techniques}

\subsubsection{Temporal Network
Evolution}\label{temporal-network-evolution}

Analysis of how terminology networks evolve over time reveals conceptual
shifts:

\begin{equation}\label{eq:temporal_network_evolution}
\Delta G(t) = G(t+1) - G(t) = \sum_{e \in E} \delta_e(t) + \sum_{v \in V} \delta_v(t)
\end{equation}

where \(\delta_e(t)\) and \(\delta_v(t)\) represent edge and vertex
changes over time periods.

\textbf{Key Evolutionary Patterns}:

\begin{itemize}
\tightlist
\item
  \textbf{Network Growth}: Addition of new terms and relationships
\item
  \textbf{Structural Rearrangements}: Changes in network topology
\item
  \textbf{Conceptual Consolidation}: Strengthening of established
  relationships
\item
  \textbf{Paradigm Shifts}: Major restructuring events
\end{itemize}

\subsubsection{Multi-Scale Network
Analysis}\label{multi-scale-network-analysis}

Extending network analysis to multiple scales reveals hierarchical
organization:

\textbf{Local Scale}: Individual term relationships within domains
\textbf{Domain Scale}: Inter-term relationships within domains
\textbf{Cross-Domain Scale}: Relationships between domains \textbf{Field
Scale}: Relationships across the entire entomological terminology
network

\subsection{Extended Validation
Frameworks}\label{extended-validation-frameworks}

\subsubsection{Inter-Subjectivity
Validation}\label{inter-subjectivity-validation}

Advanced validation incorporates multiple perspectives:

\textbf{Expert Validation}: Entomological domain experts review
classifications \textbf{Peer Validation}: Interdisciplinary researchers
assess cross-domain mappings \textbf{Historical Validation}: Analysis of
terminology evolution against known conceptual shifts
\textbf{Cross-Cultural Validation}: Comparison with non-English
entomological literature

\subsubsection{Robustness Testing}\label{robustness-testing}

Robustness analysis ensures result stability:

\textbf{Subsampling Stability}: Performance across different corpus
subsets \textbf{Parameter Sensitivity}: Robustness to algorithmic
parameter variations \textbf{Annotation Consistency}: Agreement across
multiple human annotators \textbf{Temporal Stability}: Consistency
across publication periods

\subsection{Advanced Case Study
Analysis}\label{advanced-case-study-analysis}

\subsubsection{Caste Terminology Evolution:
1850-2024}\label{caste-terminology-evolution-1850-2024}

Ultra-longitudinal analysis reveals century-scale conceptual evolution:

\textbf{Pre-Darwinian Period (1850-1859)}: Essentialist caste categories
based on morphological differences

\textbf{Darwinian Synthesis (1860-1899)}: Evolutionary explanations for
caste differences

\textbf{Genetic Revolution (1900-1949)}: Chromosomal mechanisms
underlying caste determination

\textbf{Molecular Biology Era (1950-1999)}: Gene expression and hormonal
control of caste differentiation

\textbf{Genomic Era (2000-2024)}: Epigenetic and transcriptomic
regulation of caste phenotypes \cite{chandra2021epigenetics},
accompanied by calls to broaden conceptions of sociality beyond
traditional eusocial models \cite{sociable2025}. \citet{warner2024caste}
demonstrate that caste differentiation becomes increasingly
\emph{canalized} from early development through cascading
gene-expression changes modulated by juvenile hormone signaling, while
gene expression in \emph{Lasius niger} is more strongly influenced by
age than by caste---further undermining the fixedness implied by
``caste'' terminology.

\subsubsection{Superorganism Concept
Evolution}\label{superorganism-concept-evolution}

Detailed analysis of the superorganism concept across seven decades:

\begin{table}[h]
\centering
\begin{tabular}{|l|c|c|c|c|}
\hline
\textbf{Era} & \textbf{Dominant Metaphor} & \textbf{Key Evidence} & \textbf{Critiques} & \textbf{Legacy} \\
\hline
1960s & Organismic & Division of labor analogies & Ignores individual variation & Established field \\
1970s & Cybernetic & Communication networks & Mechanistic reductionism & Systems thinking \\
1980s & Genetic & Kin selection theory & Haplodiploidy focus & Evolutionary framework \\
1990s & Neuroendocrine & Pheromonal control & Colony complexity & Regulatory mechanisms \\
2000s & Epigenetic & DNA methylation & Environmental effects & Developmental plasticity \\
2010s & Microbiome & Symbiont communities & Host-symbiont dynamics & Extended organism concept \\
\hline
\end{tabular}
\caption{Evolution of superorganism concept across research eras}
\label{tab:superorganism_concept_evolution}
\end{table}

\subsection{Methodological
Reflections}\label{methodological-reflections}

\subsubsection{Mixed-Methodology
Integration}\label{mixed-methodology-integration}

Our approach successfully integrates qualitative and quantitative
methods:

\textbf{Qualitative Contributions}:

\begin{itemize}
\tightlist
\item
  Theoretical framework development
\item
  Conceptual category identification
\item
  Historical context analysis
\item
  Cross-domain relationship mapping
\end{itemize}

\textbf{Quantitative Contributions}:

\begin{itemize}
\tightlist
\item
  Statistical pattern identification
\item
  Network structure analysis
\item
  Temporal trend quantification
\item
  Validation metric development
\end{itemize}

For a discussion of methodological limitations and scope considerations,
see Section \ref{sec:discussion}. Future research directions, including
advanced semantic analysis (transformer-based embeddings, multilingual
extensions) and practical applications (terminology standards, peer
review tools), are discussed in Section \ref{sec:conclusion}.

\newpage

\section{Supplemental Applications}\label{sec:supplemental_applications}

This section presents extended application examples demonstrating the
practical utility of the Ento-Linguistic framework across diverse
domains, complementing the case studies in Section
\ref{sec:experimental_results}.

\subsection{Biological Sciences
Applications}\label{biological-sciences-applications}

\subsubsection{Evolutionary Biology}\label{evolutionary-biology}

Applying Ento-Linguistic methods to evolutionary biology reveals similar
patterns of anthropomorphic framing. Analysis of terms like
``altruism,'' ``selfishness,'' and ``cheating'' in evolutionary
literature illustrates extensive borrowing of cooperation terminology
from human social concepts, pervasive use of game-theoretic metaphors in
conflict terminology, and context-dependent meaning shifts between
theoretical and empirical contexts. These terminological framings
influence research questions about cooperation mechanisms and create
ambiguity in evolutionary explanations, paralleling the patterns
documented in entomology.

\textbf{Worked example --- Kin-selection terminology network.} Running
the \texttt{TerminologyExtractor} over 200 abstracts from
\emph{Behavioral Ecology and Sociobiology} (2010--2020) produces a term
co-occurrence graph with three prominent clusters: (1) a \emph{strategy}
cluster (``altruism,'' ``cheating,'' ``punishment,'' centering on
game-theoretic metaphors), (2) a \emph{mechanism} cluster (``gene
expression,'' ``pheromone,'' ``receptor''), and (3) a \emph{scale}
cluster (``colony,'' ``population,'' ``kin group''). The
\texttt{DomainAnalyzer} identifies 47 cross-cluster edges, 31 of which
involve anthropomorphic framing --- a result comparable to the 62\%
anthropomorphic edge rate found in the core entomology corpus. The
\texttt{ConceptualMapper.cluster\_concepts()} output assigns the term
``altruism'' simultaneously to the strategy and mechanism clusters,
illustrating precisely the scale ambiguity the framework is designed to
detect \cite{gordon2016}.

The pipeline invocation follows the standard orchestration pattern:

\begin{verbatim}
from src.analysis.term_extraction import TerminologyExtractor
from src.analysis.domain_analysis import DomainAnalyzer

extractor = TerminologyExtractor()
terms = extractor.extract_terms(abstracts, min_frequency=3)
# terms["altruism"].domains → ["Behavior and Identity", "Economics"]
# terms["altruism"].confidence → 0.87

analyzer = DomainAnalyzer()
results = analyzer.analyze_all_domains(terms, abstracts)
# results["Behavior and Identity"].ambiguity_metrics["mean_ambiguity"] → 0.73
\end{verbatim}

\subsubsection{Ecology}\label{ecology}

Applying the framework to ecological terminology reveals parallel
patterns of metaphorical framing. Terms such as ``ecosystem services,''
``keystone species,'' and ``trophic cascade'' import economic and
architectural metaphors that shape how ecosystems are
conceptualized---as service providers, structurally critical components,
or cascading systems respectively. Running the \texttt{DomainAnalyzer}
over a corpus of 100 conservation biology abstracts reveals that 58\% of
key terms carry economic framing (cost--benefit assumptions), while
``keystone'' imposes an architectural hierarchy that may obscure the
distributed redundancy characteristic of many ecological networks. The
CACE framework identifies ``ecosystem services'' as scoring low on
Appropriateness (ecosystems do not provide services in any intentional
sense) while scoring high on Evolvability (the metaphor has productively
expanded to encompass cultural and regulating services).

\subsubsection{Neuroscience}\label{neuroscience}

Ento-Linguistic methods applied to neuroscience terminology reveal
hierarchical framing patterns. Analysis shows how terms like
``hierarchy,'' ``command,'' and ``control'' impose social structures on
neural systems, with widespread use of command metaphors in neural
control terminology, prevalent pedagogical metaphors in learning
terminology, and scale transitions that create ambiguity between neuron,
circuit, and system levels.

\textbf{Worked example --- Motor-control terminology.} Applying the
\texttt{DiscourseAnalyzer.analyze\_discourse\_patterns()} method to a
curated corpus of 50 motor-neuroscience review articles detects
\emph{hierarchical\_framing} in 82\% of texts, primarily through the
terms ``command neuron,'' ``executive control,'' and ``motor program.''
The \texttt{quantify\_framing\_effects()} method further reveals that
texts using command metaphors also exhibit higher rates of teleological
language (framing\_strength = 0.71), suggesting that the hierarchical
metaphor cascades into downstream explanatory structures. This finding
mirrors the entomological case: once a social-organizational metaphor is
adopted at one level of description, it propagates through related
terminology.

\subsection{Historical and Cross-Cultural
Analysis}\label{historical-and-cross-cultural-analysis}

\subsubsection{Longitudinal Terminology
Studies}\label{longitudinal-terminology-studies}

Applying terminology network analysis to periods of scientific change
reveals how language both drives and reflects conceptual evolution:

\begin{itemize}
\tightlist
\item
  \textbf{Darwinian Revolution (1830--1870)}: Shift from creationist to
  naturalistic explanatory frameworks
\item
  \textbf{Molecular Biology Revolution (1940--1970)}: Transition from
  classical to molecular explanations
\item
  \textbf{Genomic Era (2000--present)}: The rise of ``-omics''
  terminology and its effects on conceptual framing
\end{itemize}

Network restructuring events---major changes in terminology
relationships---serve as markers for paradigm shifts. Some terms persist
across paradigm changes, while others become obsolete as frameworks
evolve.

\subsubsection{Multilingual Scientific
Terminology}\label{multilingual-scientific-terminology}

Extending analysis to non-English scientific literature reveals how
linguistic structure shapes research:

\begin{itemize}
\tightlist
\item
  \textbf{German}: Comparing \emph{Staaten} (``states'') vs.~English
  ``colony'' reveals fundamentally different conceptual framings of
  social insect organization
\item
  \textbf{French}: Analysis of hierarchical vs.~egalitarian conceptual
  frameworks in biological descriptions
\item
  \textbf{Chinese}: Examining how traditional concepts influence modern
  scientific language
\end{itemize}

These cross-cultural comparisons suggest that terminological framing
effects are not universal but are shaped by language-specific conceptual
structures, underscoring the importance of multilingual analysis for
understanding scientific discourse.

\subsection{Tools, Education, and
Standards}\label{tools-education-and-standards}

\subsubsection{Research Tools}\label{research-tools}

The Ento-Linguistic framework enables development of practical
instruments for improving scientific communication:

\begin{itemize}
\tightlist
\item
  \textbf{Terminology analysis software} for automated identification of
  framing assumptions in scientific texts
\item
  \textbf{Writing assistance tools} providing real-time feedback on
  terminological clarity and appropriateness
\item
  \textbf{Peer review frameworks} integrating language analysis to
  improve manuscript quality
\end{itemize}

\subsubsection{Educational Applications}\label{educational-applications}

Ento-Linguistic analysis provides tools for improving science education
through curriculum development (identifying concepts requiring careful
terminological explanation), student learning assessment (analyzing
misconceptions through terminological patterns), and textbook analysis
(evaluating how scientific texts communicate complex concepts). Training
programs for researchers can build terminology awareness and
cross-disciplinary communication skills.

\subsubsection{Policy and Ethics}\label{policy-and-ethics}

Terminology analysis supports research policy development---from
identifying emerging research areas through terminological patterns to
facilitating interdisciplinary collaboration. Ethical applications
include promoting inclusive language that avoids cultural bias, ensuring
transparent communication that serves research goals, and developing
responsible guidelines for scientific naming practices
\cite{betternamesproject2024}.

\subsubsection{Decolonizing Entomological
Curricula}\label{decolonizing-entomological-curricula}

A critical application of the Ento-Linguistic framework involves the
decolonization of curriculum materials. Our analysis of the Power \&
Labor domain reveals that standard textbook descriptions of ant colonies
frequently rely on ``settler science'' metaphors---conquest, slavery,
and colonial expansion---that were explicitly cultivated during the
imperial era to naturalize colonial projects
\cite{mavhunga2018transient}.

\textbf{Curriculum Audit Protocol}: We propose a
\texttt{CurriculumAuditor} module that scans educational texts for three
specific colonial narrative tropes:

\begin{enumerate}
\def\labelenumi{\arabic{enumi}.}
\tightlist
\item
  \textbf{The Civilizing Mission}: Framing ``advanced'' eusocial insects
  as superior to ``primitive'' solitary species, mirroring colonial
  development narratives.
\item
  \textbf{The Frontier Myth}: Describing territory expansion as
  ``manifest destiny'' or ``empty land'' colonization, ignoring
  competitive exclusion or incumbent species.
\item
  \textbf{The Plantation Model}: Describing fungus-farming ants solely
  through the lens of industrial agriculture and labor management,
  obscuring symbiotic complexity.
\end{enumerate}

By identifying these tropes, educators can reframe lessons to emphasize
ecological integration, symbiosis, and diverse social strategies, moving
away from narratives that implicitly validate colonial ideologies
\cite{laciny2024terminology}.

\subsection{Future Directions}\label{future-directions-2}

Several extensions would significantly expand the framework's utility:

\textbf{Machine learning classification} of framing types could automate
the detection of anthropomorphic, hierarchical, and economic framings at
scale. \textbf{Advanced network analysis} using temporal graph methods
could track terminology evolution in real time. \textbf{Ontology
integration}---mapping to existing biological ontologies---would ground
the framework in established knowledge structures.

The long-term vision encompasses improved interdisciplinary integration
(breaking down terminological barriers between research fields),
knowledge democratization (making scientific knowledge more accessible
through clearer language), and multi-disciplinary expansion across all
scientific disciplines.

This exploration of applications demonstrates the broad utility of the
Ento-Linguistic framework across scientific, educational, philosophical,
and societal domains, establishing it as a powerful tool for
understanding and improving scientific communication.

\newpage

\bibliography{/Users/4d/Documents/GitHub/template/projects/ento_linguistics/manuscript/references.bib}

\end{document}
